\quellseite{221}

\part{Das vierte Buch Mose}

\chapter{Das 1. Kapitel}

\section{Vorrede}

Dies vierte Buch hat den Namen in der lateinischen Sprache, daß es heißt das Buch der Zahl. Denn nachdem die Hütten oder Kirche zugerichtet, und aller Gottesdienst angerichtet und verordnet war, befahl Gott, man sollte alle Stämme zählen, V. 3., und einem jeden seinen Ort im Lager bestimmen, auf daß die Wohnung Gottes mitten unter ihnen wäre, und sonderlich, daß die Leviten müßten, wenn man reiste, was jedermann zu verwahren hätte; solches ist im 1., 2., 3. u. 4. Kapitel verzeichnet.

\section{Zahl der streitbaren Männer in Israel}
\quellseite{222}

\chapter{Das 2. Capitel}

\section{Ordnung der Lager im Volk Israel}
\quellseite{223}

\chapter{Das 3. Capitel}

\section[Zählung der Leviten und ihr Amt. Der Erstgebornen]{Zählung der Leviten und ihr Amt. Der Erstgebornen Auslösung}
% No apparatus blocks retained on this page. See selection.json for block-by-block exclusion reasons.
\quellseite{225}

\chapter{Das 4. Kapitel}

\section{Besondere Amtsverwaltung und Zahl der Leviten}

a Gelbe Decke) hieraus scheint, daß die innersten Teppiche sind nicht bunt, sondern eine jegliche hat ihre eigene Farbe allein gehabt, welche gelb ist, die ist ganz gelb gewesen, und welche weiß, ganz weiß etc.
% No apparatus blocks retained on this page. See selection.json for block-by-block exclusion reasons.
\quellseite{227}

\chapter{Das 5. Kapitel}

\section{Reinigung des Lagers. Versöhn- und Eiferopfer}
\quellseite{228}

\section{Summarium}

Dies Kapitel ordnet, wenn einer den andern beleidiget hat, wie man den Schaden widerlegen und den fünften Theil über die Hauptsumma soll geben, auf daß der, so beleidiget ist, versühnet werde.

Darnach ist eine Ordnung, wie sich ein Mann halten soll gegen seinem Weibe, wenn sie Unzucht halben verdächt ist, V. 11. Und zeigt Gott hie sonderlich an, wie er dem Ehebruch feind sey, und ihn mit Unfruchtbarkeit des Leibes auch andern hartem Unglück strafen wolle.

\section{Votum}

Helfe der treue barmherzige GOtt, daß wir alle seine Ordnung allezeit fleißig halten, niemand Schaden thun, und daß wir Junge und Alte, in Zucht, auch Keuschheit, in und außer dem Ehestande leben, Amen.

\chapter{Das 6. Kapitel}

\section{Vorrede}

Undank ist der Menschen Untergang. Nun müssen wir alle bekennen, daß die Wohlthaten, so uns Gott aus Gnaden und Barmherzigkeit erzeiget hat und noch täglich erzeiget, so wenig können gezählet werden, als der Menschen Haar auf dem Haupt, als die Sterne am Himmel und als die Sandkörnlein am Ufer der Wasser. Wer wollte denn nicht Gott danken, loben und preisen täglich für solche seine unzählige, große und nützliche Wohlthaten? Danket demnach dem HErrn eurem Gott von ganzem Herzen, und ehret seinen Namen ewiglich, Ps. 86. Dieses vermahnet auch Sirach im 40. Kapitel, da er also saget: Singet löblich, und lobet den HErrn in allen seinen Werken, preiset seinen Namen herrlich, danket ihm und lobet ihn mit Singen und Klingen, und sprechet also im Danken: Alle Werke des HErrn sind sehr gut, und was er gebeut, das geschiehet zu rechter Zeit. Darauf höret ferner an den ersten Theil des 6. Kapitels.

\section[Gesetz der Nazaräer oder Verlobten. Formular das Volk]{Gesetz der Nazaräer oder Verlobten. Formular das Volk zu segnen}

a Auf Hebräisch heißt dies Neser, und wer sie hält, heißt Nasir, welchem nach auch unser HErr Jesus Christus Nazarenus heißt, und er der rechte Nasir ist.
\quellseite{229}

\section{Votum}

Wir danken GOtt herzlich für alle seine vielfältigen und unaussprechlichen Wohlthaten; der wolle uns ferner an Leib und Seele, zeitlich und ewiglich, segnen und behüten, um JEsu Christi, seines lieben Sohns und unsers HErrn willen, Amen.

\section{Der andere Theil des 6. Capitels}

\section{Vorrede}

Der Ursprung und Quellbrunn aller Sünden ist, wenn man Gottes Wort hinterrück wirft; wie dagegen aller guten Werke Ursprung ist, wenn man Gott Gehorsam leistet, und sein Wort Norm und Form sein läßt, darnach wir unser Leben anstellen: also, die erste Weihe zum Verderben ist, wenn man die Zucht hasset, ungestraft sein will und Gottes Wort verwirft. Wenn man aber die Sünde erkennet und sich den Geist Gottes willig fragen läßt, so ist es der erste Antritt zur Seligkeit. Derwegen danken wir heute abermal Gott billig und herzlich, daß er uns sein heilig Wort hören läßt, und bitten, er wolle uns durch den heil. Geist erleuchten und helfen, daß wir dasselbe zu täglicher Besserung und Erquickung fruchtbarlich gebrauchen mögen, Amen. Darauf höret an den andern Theil des 6. Capitels.
\quellseite{230}

\section{Summarium}

Ist eine Ordnung, wenn das Volk beisammen ist, mit was Worten der Priester sie segnen soll, V. 22. Denn wie man siehet, Gott will, daß man bei seinem Dienst thue, daß man es alles nach seinem Wort und Befehl thun und aus eigenem Gutdünken nichts fürnehmen soll.

Nun fasset aber dieser Segen drei Stücke: Daß er am ersten sagt, V. 24.: Der HErr segne dich und behüte dich, daß gehet auf den leiblichen Segen, daß Gott Leib und Leben, Nahrung und alles segnen und darnach für aller Gefahr und Schaden, so der Teufel und die arge Welt zufügen kann, behüten wolle. Denn es liegt nicht an unserer Arbeit noch Fleiß, wiewohl es beides sein soll, sondern an Gottes Segen und Behüten.

Das andere: Der HErr lasse sein Angesicht leuchten über dir, und sei dir gnädig, V. 25. Das gehet auf den geistlichen Segen, daß Gott gnädig sein und sein Wort rein und lauter geben wolle, auf daß wir ihn recht erkennen und nach seinem Wort uns halten mögen.

Das dritte: Der HErr erhebe sein Angesicht über dich, und gebe dir Frieden, V. 26. Dies gehet auf den geistlichen Schutz; denn wo Gottes Wort und Geist ist, da feiert der Teufel nicht, sicht die Leute immer an, ob er sie davon möchte abführen. Da ist denn vonnöthen, daß Gott sein Angesicht über uns hebe, das ist, daß er auf uns schaue und in solchen Anfechtungen erhalte, Friede und Sieg gebe wider den Teufel und alles Unglück.

Dabei sollen wir auch lernen, wo solcher leiblicher und geistlicher Segen und Schutz ist, daß man Gott von Herzen dafür danken soll und erkennen, daß wirs nicht von uns selbst, sondern von Gott haben. Darnach, daß man auch täglich bitten soll, daß Gott solche Gaben uns erhalten und mehren wolle.

\section{Votum}

Wir danken GOtt herzlich für alle seine Wohlthaten; der wolle uns ferner geben an Leib und Seele, was uns gut und nöthig ist, und wolle uns auch behüten für allem, was uns an Leib und Seele böse und schädlich ist, Amen.

\chapter{Das 7. Kapitel}

\section{Vorrede}

Eine jegliche Obrigkeit, vom höchsten Amt bis zum geringsten auf Erden, soll sich theilen als in drei Orte, und soll sehen, erst über sich, darnach in sich, endlich neben sich. Erstlich soll die Obrigkeit sehen über sich, auf den, der da sagt in Sprüchen Salomonis am 8. Kapitel: Mein ist Rath und That; und ist der Erb-Herr über alle Heiden, Ps. 82. Mit desselben Hülfe und Segen sollen sie ihr Amt anfangen, um den Schlüssel zur Regierung bitten, und Gottes Ehre, Wort und Kirche helfen, fördern und schützen. Zum andern soll die Obrigkeit sehen und gehen in sich selbst, das ist, in ihr Gewissen, und bedenken, wie sie regieret, was sie ordnet und wie sie ihr Amt ausführet, daß sie das Gewissen zum Zeugen habe, recht und fleißig geführten Regiments. Zum dritten soll auch die Obrigkeit neben sich sehen, das ist, eine jede auf ihre befohlenen Unterthanen, daß sie gegen den Frommen Friedrich und gegen den Bösen Ernst sei. Dies erinnert auch an einem Stück ein Theil des 7. Kapitels.

\section[Opfer und Geschenke der Fürsten zur Einweihung der]{Opfer und Geschenke der Fürsten zur Einweihung der Stiftshütte}
\quellseite{231}

a) Löffel) oder ein Köpflein oder sonst inwendig rund wie die Salbkrüglein.
% No apparatus blocks retained on this page. See selection.json for block-by-block exclusion reasons.
\quellseite{233}

\section{Summarium}

Ist ein fein Exempel, wie die zwölf Fürsten den Gottesdienst anheben, und die Ersten sind, so mit ihrem Opfer und reichlichem Geschenk zum Gottesdienst helfen. Denn gleich wie die Obrigkeit sonst in allen andern Sachen ein gut Exempel geben, und aller Tugend und Ehrbarkeit sich fleißigen soll: also ist's sonderlich der Obrigkeit Amt, sich Gottes Wort und des rechten Gottesdienst mit Ernst annehmen, und den Unterthanen ein gut Exempel fürtragen. Bei solchen Regenten ist sich alles Glücks und Heils zu versehen.

Widerum wo gottlose Obrigkeit ist, die entweder Gottes Wort nicht achtet, wie gemeiniglich geschiehet, oder auch wohl verfolget, da wird alles Unglück sich zeitlich genug finden, wie die Exempel vor Augen sind.

\section{Votum}

Wir danken Gott herzlich für unsere liebe Obrigkeit, unter deren Schutz mir Gottes Wort und den rechten Gottesdienst haben, der wolle sie gnädiglich segnen, schützen und erhalten, zu Lob und Preis seines Namens, und zur Wohlfahrt der Unterthanen, Amen.

\chapter{Das 8. Kapitel}

\section{Einweihung der Leviten}

a Vorwärts) das ist, die Seiten oder Schnauzen an den Lampen soll er gegen dem Vorhang richten, daß es für dem Leuchter und Vorhang Licht sei.

b Sünd-Wasser) entsündigen ist so viel als abspülen oder lossprechen; daher das Wasser, damit sie abgespület wurden, heißt Sünd-Wasser.
\quellseite{234}

\section{Summarium}

Das achte Kapitel hält in sich die Ordnung mit der Leviten Weihe oder Heiligung. Wie aber solchem Exempel nach der Papst mit seinen Pfaffen gehandelt hat, ist für Augen. Denn seine Pfaffen müssen alle geweyhet und gesalbet sein. So doch solches weder Christus befohlen, noch die Apostel geordnet haben. Es ist aber kein gerathen. Des HErrn Christi und der Apostel Pfaffen hatten eine Weihe, die heißt die Weihe des heil. Geistes, dieselbe hat der Papst und seine Pfaffen nicht. Darum sollen sie ohne Gottes Befehl ihnen eine eigene erdenken, die doch sehr schlecht ist, und ein geringes Ansehen auch für der Welt hat. Denn was soll es zum Amt helfen, daß man sie mit Oele schmieret. In Summa, der Papst als ein Affe, hat ohne Befehl das hiernach gethan, welches die Juden aus Befehl Gottes gethan haben, daß sie keinen zum Gottesdienst haben genommen unter fünf und zwanzig Jahren, und daß sie im fünfzigsten Jahr sind frei gewesen; wiewohl der Papst solchs mit der Freiheit auch geendet hat.

\chapter{Das 9. Kapitel}

\section{Vom Paßah, und Zeichen zum Lagern oder Fortziehen}
\quellseite{235}

\section{Summarium}

Im neunten Kapitel ist nichts Sonderliches zu merken, denn wie Moses preiset, daß die Kinder Israel mit ihrem Ziehen und Stillliegen sich nach Gottes Wort und Befehl gehalten haben. Solches sollten wir lernen, daß es die größte Kunst und Tugend ist, davon der Papst mit seiner Kirchen nichts weiß, sintemal sie alle ihre Gottesdienste aus eigenem Gutdünken ohne Gottes Befehl und Wort haben angericht.

\chapter{Das 10. Kapitel}

\section{Gebrauch der Trommeten. Der Israeliten Fortreisen}
\quellseite{236}

a Richteten) indes die Kahathiter das Heiligthum hernach trugen, waren hierbei vorhin die Gersoniter und Merariter, und richteten die Wohnung auf, daß die Kahathiten mit der Lade die Wohnung bereit fanden.

\section{Summarium}

Wie ist sonderlich zu merken das Gebet Mos., V. 35, welches allen Regenten zum Exempel ist vorgeschrieben, daß sie neben allem andern, das sie ihres Amts halben thun sollen, auch beten und bitten sollen, auf daß Gott seine Feinde stürzen und die Seinen gnädiglich beschützen und segnen wolle.

Denn obwoll die Obrigkeit mit weltlicher Gewalt die Bösen strafen und die Frommen beschützen kann, so dürfen sie doch fürnehmlich der Hülfe Gottes dazu, daß es wohl von Statten gehe, und der Teufel solches nicht hindere, wie er wohl im Sinne hat, und zu thun sich unterstehet, und ihm bei dem meisten Theil geräth, daß sie ihres Amts vergessen, die Sünde nicht strafen, wie es Gott befohlen hat, und die Tugend und die Ehrbarkeit nicht schützen noch handhaben.

\section{Votum}

Verleihe GOtt, der HErr aller Herren, daß die Obrigkeit im Gebet und allen andern Stücken ihres Amts sich recht und christlich verhalte, damit alle fromme Unterthanen ihm dafür Lob und Dank sagen mögen allezeit, immer und ewiglich, Amen.
\quellseite{237}

\chapter{Das 11. Kapitel}

\section{Vorrede}

Der heilige Chrysostomus lehret von Kraft des gläubigen Gebets also: Solche große Herrlichkeit und Wirkungen hat das Gebet in sich, daß es aus Menschen Gottes Tempel macht. Und wie Gold, Silber, Edelsteine und Marmor der Könige Häuser schön und herrlich machen: also schmückt und zieret das Gebet einen Christen, und macht aus ihm einen schönen, herrlichen, gezierten Tempel und Palast Christi, darin er wohnet; und welchen der Himmel nicht umgreift, der gibt sich in die Seele, in welcher fleißig und emsig Gebet und Seufzen ist. Ferner sagt auch Chrysostomus: Viel Heilige oder Gläubige Gottes haben durch die Macht oder Kraft ihres Gebets erhalten nicht alleine Land und Leute, sondern auch die ganze Welt. Derwegen laßt uns auch fleißig und in wahrem Glauben alle Tage beten. Solches unsers herzlichen Gebets werden wir gewiß genießen, zu unsrer geistlichen und leiblichen Wohlfahrt, um Jesu Christi unsers Herrn Verdienst und Vorbitt willen. Darauf höret nun an das 11. Kapitel.

\section{Das lüsterne Volk wird gestraft}

\section{Summarium}

Dies ist ein schrecklich Exempel, wie die Juden des Himmelbrods in der Wüsten überdrüssig worden sind, und wünschen um des Fleisches und andere Essens willen, daß sie wieder in Egypten wären, haben vergessen aller Noth und Jammers, den sie darin erlitten haben, und Gott sie davon erlöset, V. 1. Solche Unbarkeit straft Gott mit Feuer vom Himmel. Darum sollen wir lernen, daß wir um zeitlichen Mangels willen nicht auch der herrlichen und geistlichen Güter vergessen, wie jetzo sehr gemein ist. Denn daß wir der Sünden los und zur rechten Erkenntniß Gottes und wahren Gottesdienst kommen sind, solches halten viel Leute für ein Geringes und Schlechtes; und klagen dieweil über den zeitlichen Mangel, daß
\quellseite{238}
mehr Mißwachsung der Früchte, Krankheit, Krieg und anders ist, denn unter dem Papstthum, so doch solches nicht des Evangelii, sondern unsere Sünde, Unbußfertigkeit und greuliche Undankbarkeit schuldig ist.

Es ist auch ein fein Exempel, wie Gott Mosen und des Volks Gebet erhöret, und mit der Strafe aufhöret.

\section{Summarium}

Ist ein fein Exempel, wie Gott sich bewilliget dem Volk, ihrem Begehr nach, Fleisch zu geben. Aber Moses achtet, es sei unmöglich. Denn er rechnet nur den Haufen, und denket nicht, was Gott vermöge zu thun. Darum sollen wir an solchem Exempel lernen, wie leicht es Gott ist, uns zu geben und zu thun, was uns gleich unmöglich dünket. Denn Gottes Hand ist nicht verkürzt, das ist, Gott ist allmächtig; ob bei uns gleich etwas unmöglich ist, sollen wir auf seine Gewalt und Allmächtigkeit sehen, und uns derselben in allerlei Noth und Mangel trösten.

Gott strafet hie zum andernmal die Juden, um ihrer Undankbarkeit und Murrens willen wider Mosen und den HErrn.

\section{Votum}

GOtt, der HErr, behüte uns für Undankbarkeit, Unglauben, Ungeduld, und allen andern Sünden, und verleihe uns allen Gnade, daß wir täglich fleißig und fruchtbarlich beten, und christlich leben, Amen-

13*
\quellseite{239}

\chapter{Das 12. Kapitel}

\section{Vorrede}

Gott ist tausendmal geneigter die Sünde zu vergeben, denn wir darum zu bitten, und siehet nicht darauf, wie groß die Sünden, und wie viel derselben sind. Ja er wendet auch ab, oder lindert gnädiglich die mit Sünden verdieneten Strafen. Hieher gehört, daß Gott selber saget durch den Propheten Esaiam im 65. Capitel: Es soll geschehen, ehe sie rufen will ich antworten, wenn sie noch reden, will ich hören. Und im 1. Kapitel spricht Esaias: Wenn eure Sünde gleich blutroth ist, soll sie doch schneeweiß werden, und wenn sie gleich ist wie Rosinfarbe, soll sie doch wie Wolle werden. Deß sollen sich trösten alle bußfertige Sünder und betrübte Kreuzträger; sie sollen auch Buße thun, und also Gottes väterlicher, gewisser Hülfe in beständiger Geduld erwarten. Davon höret an das 12. Kapitel.

\section{Die murrende Mirjam wird aussätzig}

a (Geplagter) elender, der viel leiden mußte, Ps. 132, 1. Gedenke Davids und alles seines Leidens, Ps. 18, 36. Ps. Sal. 18, 12. Ante gloriam passio. Das ist, vor der Herrlichkeit Leiden.

\section{Summarium}

Ist eine Historia, wie Aaron und Mirjam Mose Schwester sich ihres Standes und Wesens überheben, und Mosen verachten, V. 2. Solche Hoffart gefiel Gott übel, darum straft er sie beide mit harten Worten und Mirjam mit dem Aussatze, auf daß wir lernen und demüthigen, und niemand verachten. Sonderlich aber, daß wir die in Ehren halten, so Gottes Amt führen und mit dem Wort umgehen. Darnach ist ein fein Exempel, wie Aaron seine Sünde erkennet, und Mosen bittet, daß er gegen Gott solche Sünde abbitten und seiner Schwester wieder vom Aussatz helfen wolle, V. 12. Denn das gehöret darauf, wer unrecht gethan hat, daß ers ihm lasse leid sein, nimmer thue, und um Gnade bitte, so will Gott die Strafe auch gnädig abwenden, wie man hie siehet.

\section{Votum}

Der treue, barmherzige GOTT wolle uns auch gnädig sein, und nicht strafen, wie wirs alle verdienet haben, und wolle uns regieren, daß wir ihn lieben, und alle die ehren, denen wir Ehre schuldig sein, Amen.

\chapter{Das 13. Kapitel}

\section{Vorrede}

Gott ist allmächtig, der kann alles thun, und mehr, denn wir in unsern Nöthen glauben, und
\quellseite{240}
ihm trauen können. So ist er auch wahrhaftig, was er zusagt, das hält er fest und erfüllets endlich. Denn von Gottes Allmacht sagt David im 135 Psalm: Alles was der HErr will, das thut er, im Himmel, auf Erden, im Meer und in allen Tiefen. Und St. Paulus Ephes. 3.: Gott kann überschwänglich thun, über alles das wir bitten oder verstehen. Von Gottes Wahrheit aber spricht auch David Psalm 89.: Gnade und Wahrheit sind für des HErrn Angesichte. Und St. Paulus Röm. 3.: Gott ist wahrhaftig und alle Menschen falsch. Derwegen wollen wir in diesen gefährlichen Zeiten Gott vertrauen, und ihm uns und die Unsern befehlen, er wird uns wohl schützen und erhalten, und sollt es allen Teufeln in der Hölle leid sein. Dies unterrichtet auch das 13. Kapitel.

\section{Von Kundschaftern des Landes Canaan}

a Escol heißt eine Traube, daraus wird der Traubenbach.
\quellseite{241}

\section{Summarium}

Ist eine Historia, wie Moses zwölf Männer ausschickt, das Land Canaan zu besichtigen, V. 4., und dieselbigen das Volk verzagt machen, daß sie daran zweifelten und besorgten, sie würden den Leuten zu schwach sein, und das Land nicht mögen einnehmen, V. 32. Solches ist ein fein Exempel, in welchem man siehet, wie die Vernunft, so bald sie von der Zusagung Gottes abfällt und sich nicht daran hält, daß Gott allmächtig und gnädig sei, an ihr selbst verzagen und verzweifeln muß.

Darum sollen wir uns lernen trösten der Hülfe und Verheißung Gottes, es stehe unserthalben wie es wolle, daß doch unser Vertrauen und Trost allein auf Gottes Wort, Güte und Allmächtigkeit stehe; so soll es nicht Noth haben.

\section{Votum}

Der allmächtige und treue GOtt stärke unsern schwachen Glauben, und helfe, daß wir ihm in Noth und Tod beständig trauen, und zeitlich und ewig an Leib und Seele erhalten werden, Amen.

\chapter{Das 14. Kapitel}

\section{Vorrede}

Es ist uns nichts auf Erden und in dieser Welt, (was das äußerliche Leben betrifft) höher und größer vonnöthen, denn die wahre brüderliche Liebe und Treue, als deren wir so wenig gerathen mögen, als der natürlichen Sonne. Denn, gleich wie wir ohne die Sonne nicht mögen leben, also müßten wir auch mit einander zu Grunde gehen, wo die Liebe und Treue aufhörete, und dagegen alle Untreue und Ungerechtigkeit überhand nähme. Diese Liebe sollen wir an unsern Nächsten üben, fürnehmlich aber mit dem Gebet für ihn in seiner Gefahr und Noth.

Von diesem der h. Chrysostomus also saget: daß einer für sich bittet, dazu treibet ihn seine eigene Noth, daß er aber für seinen Nächsten bittet, dazu zwinget ihn die Liebe. Es ist aber für Gott ein lieblicher Gebet, dazu uns die brüderliche Liebe, denn dazu uns die Noth treibet. Dies soll auch erinnern das 14. Kapitel.

\section{Das murrende Volk wird geschlagen}
\quellseite{242}

\section{Summarium}

Dies ist eine feine Historia, in welcher man siehet erstlich, wie bald der Pöbel zu bewegen ist, daß er von Gott und seinem Wort abfällt. Sintemal die Juden aller Wunderwerke vergessen, so Gott bei ihnen gethan hatte, als könnte oder wollte ihnen Gott nimmer helfen, und begehren wieder in Egypten, V. 4.

Zum andern ist hie eine schöne Vermahnung wider solches Verzagen und Unglauben, daß man den Gedanken nicht Raum geben, sondern mit Ernst widerstehen soll, die uns zum Unglauben treiben, wie die Juden hie, die da zaghaftig werden, daß sie gedenken, die Feinde sind ihnen zu stark und mächtig, V. 11.; sondern daran sollen wir gedenken in aller Noth, daß Gott bei uns sein, und mit uns wider unsere Feinde streiten will. Solches sollte uns mehr trösten, denn alle Welt mit alle ihrem Vermögen uns schrecken kann. Aber wer also auf Gott und seine Güte sich und andere tröstet, der muß gewarten, daß man ihn steinige, wie die Juden hie Mose zu thun gedachten, V. 13. Aber Gott wachet über die Seinen und errettet sie.
\quellseite{243}

\section{Summarium}

Ist ein fein Gebet, daß Gott schonen, und das Volk um der Sünden willen nicht gar vertilgen wolle. Dies Gebet soll uns eine Form sein, wie wir unsere Fürbitte gegen Gott sollen stellen, auf daß man in aller Fürbitte und allem Gebet erstlich Gottes Ehr und Namen bedenke. Zum andern, daß man auf Gottes Verheißung sehe, und sich seiner Barmherzigkeit tröste, weil er als ein gnädiger Gott will gerühmet sein, daß er solche Gnade gegen uns beweisen wolle. Solch Gebet kann nicht fehlen. Darum saget Gott Mose zu, er wolle vergeben. Aber gleichwohl dräuet er, und zur Warnung: es soll ihr keiner ins Land Canaan kommen, und die Kinder, so solcher Sünde unschuldig waren, sollen der Väter Schuld mit tragen vierzig Jahre; auf daß wir ja sehen, wie Gott der Sünde feind sei, und lernen uns dafür hüten, auf daß nicht unsere Kinder derselben nach uns entgelten müssen.

\section{Summarium}

Ist erstlich ein Exempel, wie Gott die Kundschafter, so zu solcher Sünde dem Volk hatten Ursach gegeben, mit dem Tode strafet. Caleb aber und Josua blieben lebendig, warum daß sie solcher Sünde gewehret, und auf Gottes Hülfe das Volk vertröstet hatten. Denn Gott will beides thun, böse Buben strafen, und die Frommen segnen und erretten.

Zum andern ist noch ein Exempel hie, in welchem wir sehen, wie die menschliche Natur so gar Gott und seinem Willen entgegen ist. Da die Juden Gottes Hülfe am meisten sich sollten getröstet haben, und ins Land Canaan ziehen, verzagten sie, und wollten mit nichten hinein. Jetzt, so es ihnen Gott verbeut und saget, sie sollen vierzig Jahre in der Wüste bleiben, wollen sie auch nicht, sondern wollen wider Gottes Willen ins Land und kriegen, V. 40. Aber weil sie es wider Gottes Befehl und Willen fürnehmen, werden sie darüber geschlagen, V. 45. Auf daß wir lernen, Gott und seinem Worte folgen, sonst soll weder Glück noch Heil bei uns sein.

\section{Votum}

GOtt, unser lieber Vater, um JEsu Christi, unsers Mittlers, willen, durch den heiligen Geist, helfe, daß wir uns an andern Leuten spiegeln, Buße thun, und Gottes Wort allzeit gehorsam sein mit Thun und Lassen, Amen.

\chapter{Das 15. Kapitel}

\section{Vorrede}

Es ist nichts leichter, als die Sünde, und nichts schwerer als die Strafen, so auf die Sünde den Unbußfertigen folgen. Derwegen wer gesündiget hat, der thue Buße, und hüte sich ferner für willentlichen muthwilligen Sünden. Denn auf viele und große muthwillige Sünden folgen viele und große verderbliche Strafen. Denn um die Sünden es so ein arg Ding ist, daß dieselbigen von der Erde gen Himmel schreien, und für Gottes Angesicht gleich treten, und die Rache und Strafe erfordern, daß er endlich seinen Zorn darauf muß ergehen lassen, wo man sich nicht durch rechte Buße zu ihm bekehret. Darauf höret auch an das 15. Kapitel.
\quellseite{244}

\section{Entheiligung des Sabbaths wird gestraft}
\quellseite{245}

a Frevel) das ist die Sünde, so nicht will Sünde sein, sondern recht haben, wie die Ketzer- und Rotten-Sünde, welche Gottes Wort und Gesetze ändert, und nach ihrem Willen deutet.

\section{Summarium}

Ist am ersten ein ernstes Gebot, daß Gott die Sünde, so aus Unwissenheit geschehen, will versühnen lassen; aber muthwillige Sünden, so fürstlich geschehen, will er nicht ungestraft lassen und sagt, V. 22: Wer solches thue, der habe den HErrn geschmähet, sein Wort verachtet, und sein Gebot zerstöret. Darum sollen solche frevele und muthwillige Sünder schlecht ausgerottet werden. Solche Predigt soll man wohl merken, und sich desto fleißiger vor Sünden hüten.

\section{Votum}

Gott, der Brunnquell aller Gnaden, erbarme dich unser, und regiere uns, daß wir nach seinen Geboten in allen Ständen christlich leben, handeln und wandeln, damit wir zeitlichen und ewigen Strafen entfliehen, Amen.

\chapter{Das 16. Kapitel}

\section{Vorrede}

Gott fürchten, der Obrigkeit und Seelsorgern gehorsam sein, und den Nächsten lieben, sollen alle, die zeitlich gesegnet und ewig erhalten sein wollen, um Christus willen aus Gnaden; dagegen, wer zeitlich und ewig Unglück, an Leib und Seele, Gut und Nahrung begehret, der sündige nur und verharre in Sünden; Herzeleid genug wird er zur Strafe haben in diesem und künftigen Leben. Weil aber die Sünde über unser Fleisch und Blut ist Meister worden, wütet und tobet derselben für und für in uns, daß ihr leicht mag eine Gelegenheit gegeben werden, so bricht sie heraus in allerlei greuliche Laster. So laßt uns täglich und fleißig beten, Gott wolle uns regieren, daß die Sünde über uns nicht herrsche, und uns in zeitlichen und ewigen Jammer nicht verführe. Darauf höret ferner an das 16. Kapitel.

\section{Gott übet Rache an der aufrührischen Rotte}
\quellseite{246}

\section{Summarium}

Ist eine Historia, wie etliche Leviten ihnen einen Anhang machten wider Aaron, und wollten auch Priester sein, und eben so viel Gewalt haben als Mose und Aaron, V. 1. Aber wie es geraten ist, wird hernach folgen; auf daß jedermann lerne, bei dem zu bleiben, dazu ihn Gott berufen hat, und desselbigen mit Fleiß auszuwarte, und nicht anders noch Höheres begehre. Denn solches ist außer dem Beruf gegangen; darum weil es wider Gottes Willen ist, kann kein Glück dabei sein, wie man siehet, daß alle Aufrührer ein bös Ende nehmen.
\quellseite{247}

\section{Summarium}

Ist eine Historia, wie Gott um der Aufrührer willen das ganze Volk wollte vertilgen, aber Moses bittets ihm ab, daß allein die Aufrührer gestraft werden, V. 22. Und hie stehet man, wie eine greuliche Sünde es ist, außer dem Beruf gehen und sich selbst ohne Beruf eindringen zum Kirchen-Amt, weil Gott Weib und Kind mit nimmt, und die Erde alles verschlingen läßt, V. 32. Solch Urtheil Gottes sollen wir fleißig merken, und uns darnach bessern lernen. Darum befiehlet Gott, man soll die Rauchpfannen aufheben zum Gedächtniß, daß jedermann bei seinem Beruf bleibe, und ohne Beruf nichts fürnehme.

\section{Summarium}

Ist ein schrecklich Exempel, in welchem man siehet, wie gar nichts andern Gottlosen hilft, weder Strafe noch andere Vermahnung. Denn Gott selbst hatte sich erzürget mit der Strafe wider die Aufrührer, daß die Erde sie sammt ihren Weibern und Kindern verschlungen hatte. Solches ist allein Gottes Werk gewesen, wie alle Vernunft muß zeugen. Aber das blind verstockte
\quellseite{248}
Volk empöret sich auf ein neues wider Mosen und Aaron, und gibt ihnen die Schuld, als hätten sie jene umbracht; wollens derhalben strafen.

Aber Gott bedeckt Mosen und Aaron mit der Wolke, daß ihnen niemand nichts thun kann, V. 42., auf daß man sehe, wie Gott die Seinen schützen und behüten wolle. Die Bösen aber strafet er, daß ihrer vierzehn tausend und sieben hundert umkommen. Aber gleichwohl ist’s tröstlich, daß Gott auch solchen Zorn ihm läßt abbitten, und der Plage durch Gebet gewehret wird, V. 48. Solche Exempel sind sonderlich wohl zu merken, auch der Ursach halben, daß man sehe, wie Gott gnädig ist, und das Gebet erhöret, und den Zorn will fallen lassen.

\section{Votum}

Der gnädige und barmherzige GOtt erhöre auch unser Gebet, regiere uns, daß wir unserer Obrigkeit und Seelsorgern Gehorsam leisten, und lindere uns die wohlverdienten Strafen, um JEsu Christi, unsers HErrn willen, Amen.

\chapter{Das 17. Capitel}

\section{Vorrede}

Fürwitzige Leute sind unglückselige Leute. Derwegen niemand andern in ihre Aemter greifen, sondern ein jeder soll seines Berufs warten, in der Furcht Gottes, und soll Gott bitten, daß er seinen Segen und glückselige Verrichtung dazu geben wolle. Also behält ein jeder ein gut Gewissen, und also wird die Arbeit seines Berufs gebenedeiet, daß sie Gott zu Ehren, ihm selber zu Nutz und dem Nächsten zur Wohlfahrt gelanget. Derwegen ein jeglicher Christ täglich in seinem Beruf betrachten soll folgende des h. Apostels Pauli Regeln und Unterricht, 1 Cor. 7.: Ein jeglicher, lieben Brüder, worinnen er berufen ist, darinnen bleibe er bei Gott. Und 1 Thess. 4. saget er: Wir ermahnen euch lieben Brüder, ringet darnach, daß ihr stille seyd und das Eure schaffet. Darauf höret auch an das 17. Kapitel.

\section{Aarons Priesterthum wird durch den grünenden Stab bestätiget}

\section{Summarium}

Ist eine feine Historia, wie Gott mit einem neuen Wunderwerk an der Ruthe Aarons bezeuget, daß das Priesteramt allein bei Aaron soll bleiben, V. 8., auf daß niemand sich ferner desselbigen anmaßen, und um solches Frevels willen umkommen sollte. Darum heißt er auch die Ruthe Aarons aufheben und bewahren, V. 10. Solches alles dienet uns dahin, daß wir uns fremder Aemter und Berufs entschlagen, und sonderlich mit Kirchendiensten ein jeder seines Berufs sich halten, und denselbigen mit Fleiß auswarten soll.

\section{Votum}

Helfe der treue GOtt, daß wir Hörer und Thäter seines Worts sein und bleiben, und segne die Arbeit unsrer Hände, zu Lob seinem heiligen Namen, und unserm Nächsten zur Freude und Wohlfahrt, Amen.
\quellseite{249}

\chapter{Das 18. Kapitel}

\section{Vorrede}

Die uns das Himmelbrod fürtragen und mittheilen, die sollen wir mit leiblichem Brod und Nahrung versorgen. Denn dieses hat Gott ernstlich befohlen, wie der h. Apostel Paulus bezeuget, 1 Cor. 9. und saget: Wisset ihr nicht, daß, die da opfern, essen vom Opfer? Und die des Altars pflegen, genießen des Altars? Also hat auch der Herr befohlen, daß, die das Evangelium verkündigen, sollen sich vom Evangelio nähren. So soll auch das darum geschehen, damit also Gottes Wort erhalten werde. Denn wenn man das heilige Evangelium verachtet und die Diener desselben Noth leiden läßt, so nimmt Gott sein Wort endlich hinweg, das ist gewiß. Wehe aber denen Leuten, zeitlich und ewiglich, die Gottes Wort nicht haben. Daneben höret an das 18. Kapitel.

\section{Amt und Unterhalt der Priester und Leviten}
\quellseite{250}

a) Bund) im Hebräischen heißt es ein Salzbund: daß, wie das Salz erhält das Fleisch unverweslich, also soll auch dieser Bund unverrücklich sein, so redet die Schrift auch 2 Chron. 13, 5. Gott hat das Reich David gegeben, und den Seinen mit einem Salzbund.

b) Fülle) das ist Most, Wein, Oel und dergleichen.

c) Besten) hebräisch heißt es das Fette; daher auch die Reichen die Fettern hießen, das ist die Besten im Volk, Ps. 17, 10. und 22, 12.

\section{Summarium}

Ist eine Ordnung der Personen, so des Kirchendiensts warten sollen, V. 5., daß das Priesteramt bei Aaron und seinen Kindern bleiben, die andern Leviten aber, ein jeder seinem Beruf nach, sonst in andern zum Kirchendienst helfen sollen, V. 6., und daß solchen Kirchenpersonen das ganze Volk Versorgung thun, und ihre Nahrung ihnen schaffen soll. Denn außerhalb des Zehnten, alles was von Vieh und Früchten des Feldes geopfert ward, an solchem hatten die Priester und Leviten auch ihren Theil.

Solches ist uns alles zum Exempel fürgeschrieben, daß jedermann mit den Seinen dazu helfen soll, auf daß die Kirchendienste recht bestellt und erhalten werden mögen; wie Christus Luc. 10. selbst spricht, da er redet mit seinen Aposteln, die er aussandte das Evangelium zu predigen: Ein Arbeiter ist seines Lohns werth. Und St. Paulus, Gal. 6.: Man soll denen das Zeitliche nicht versagen, welche das Geistliche den Leuten mittheilen.

\section{Votum}

GOtt erhalte uns sein heiliges Wort, und regiere uns, daß wir nach demselben unsern Beruf und Aemter christlich verrichten, und uns gegen seine treuen Diener allzeit auch dankbarlich erzeigen, Amen.

\chapter{Das 19. Capitel}

\section{Vorrede}

Der edelste und beste Schatz aller Menschen hie auf dieser Welt ist Vergebung der Sünden. Der edelste und beste Schatz aber in jener Welt wird sein das ewige freudenreiche Leben. Dieser Schatz werden wir theilhaftig durch den Glauben an Jesum Christum unsern Heiland, der uns solche durch sein Blut und Tod erworben hat. Dieses lehret der h. Apostel Paulus, als da er sagt, Col. 1. und Eph. 1.: An Christo haben wir (vernehmt, die Gläubigen) die Erlösung durch sein Blut, nämlich die Vergebung der Sünde. Und denn daß er spricht Röm. 6.: Die Gabe Gottes ist das ewige Leben in Christo Jesu unserm HErrn, vernehmt der uns solches verdient hat, und derer allein die Gläubigen theilhaftig werden. Demnach laßt uns Buße thun und an Christum glauben, so haben wir
\quellseite{251}
Vergebung der Sünden, und bekommen das ewige Leben, daß ist je gewißlich wahr. Dies soll auch erinnern das 19. Kapitel.

\section{Von der röthlichen Kuh und dem Sprengwasser}

\section{Summarium}

Dies Kapitel redet von einer sonderern Ceremonie, welche die Juden mit Sprengwasser von dem Sündopfer gehabt haben, sich damit zu reinigen oder zu entsündigen, V. 4. Was nun aber für Reinigkeit durch solches Wasser die
\quellseite{252}
Juden gehabt, meldet die Epistel an die Ebräer im 9. und nennet's eine leibliche Reinigkeit. Aber durch das Blut Christi, spricht er, der sich selbst ohne Wandel durch den heiligen Geist geopfert hat, werden unsere Gewissen gereiniget von den todten Werken, zu dienen dem lebendigen Gott. Darum ist es nur eine äußerliche Ceremonie gewesen, welche Gott befohlen hat. Aber die Deutung ist sehr fein gewesen: denn die rothe Kuh so zu solchem Sündopfer gehört, V. 2., hat den Sohn Gottes unsern HErrn Christum bedeutet, der ein Opfer für der ganzen Welt Sünde worden ist. Daß man aber über die Asche der verbrannten Kuh ein frisches lebendiges Wasser gegossen, und darnach die Menschen sich damit besprengt haben, V. 9., solches ist ohne Zweifel die Deutung, daß durch die heilige Taufe das Leiden unsers lieben HErrn Jesu Christi, sammt Vergebung der Sünden und ewigem Leben, ausgetheilt solle werden, wie Christus spricht, Marc. 16, 16.: Wer glaubt und getauft wird, der wird selig. Solche Taufe ist das rechte Weihwasser, das wir im neuen Testament haben, und durch jenes bei den Juden ist gedeutet worden.

Aber der Papst hat ihm ein anderes Weihwasser in seiner Kirche erdacht, damit, überredet er die Leute, soll man täglich Sünde abwaschen. Aber wo ist Gottes Wort und Befehl? In Summa, es ist nichts Gutes an allem was der Papst je und je in der Christenheit ohne Gottes Wort hat eingeführet. Man weiß aber, was für Abgötterei, Aberglaube, Zauberei und andere Gotteslästerung aus dem geweihten Wasser gefolget. Darum es billig nicht ein geweihetes, sondern ein verfluchtes und heillos Wasser ist, dadurch das Blut Christi gelästert, und die Leute schändlich betrogen sind, als fürchte sich der Teufel dafür. Wer nun wissen will wo es her komme, der hörets hie in Mose: Aber was gehet solches den Papst und uns an, was die Juden für Weise gehabt haben.

\section{Votum}

Gott, der HErr, steure allen abgöttischen Greueln, und helfe, daß wir uns des Blutes unsers HErrn Jesu Christi, und unserer Taufe mit Freuden trösten, in aller Noth und wider den Tod, im Leben und Sterben, Amen.

\chapter{Das 20. Kapitel}

\section{Vorrede}

Der Unglaube ist die größte Sünde, ja der Unglaube ist ein Quellbrunn und Ursprung aller Sünden. Denn wie die Haupttugend ist der Glaube: also ist wiederum die Hauptsünde der Unglaube. Und alle Sünden erlangen Vergebung, allein der Unglaube nicht, und verdammt keine Sünde, wo nicht der Unglaube an Christum dazu kömmt. Derwegen prüfe ein jedes unter uns sein Herz; und da es des Unglaubens wegen nicht auch wohl gestraft werden, so höre es Gottes Wort, und brauche die heiligen Sacramente, wie solche Christus geordnet; bete auch täglich: Ich glaube lieber HErr, hilf meinem Unglauben, Marc. 9., und: HErr, stärke uns den Glauben, Luc. 17. Dies erinnert auch das 20. Kapitel.

\section{Wasser aus dem Felsen. Mirjam und Aaron sterben}

a. Quia debebit semel percutere. Das ist, weil er nur einmal schlagen sollte.
\quellseite{253}

b Dubitatio est peccatum. Sed significat mysterium, quod populus legis non potest per suum Doctorem salvus fieri, et quod sub Christi tempus Moses incredulus, idest, populus, corruit. Das ist, Zweifel ist Sünde. Aber es zeigt ein Geheimniß an, daß das Volk des Gesetzes nicht kann durch seinen Lehrer selig werden, und daß zur Zeit Christi der ungläubige Moses, das ist, das Volk, dahinfallen wird.

\section{Summarium}

Ist eine Historia, wie das Volk sich wieder versündiget mit Murren wider Mosen und Aaron, darum daß sie kein Wasser hätten, und Mose und Aaron um solches Murrens willen besorgen. Gott möchte sein Wort nicht vollstrecken, und sein Wasser aus dem Fels kommen lassen, V. 3. Solchen Unglauben zeigt Gott für eine große Sünde an, und drohet Mose und Aaron, sie sollen solcher Sünde halben auch nicht in das gelobte Land kommen, V. 12. Denn Gott will das Vertrauen bei uns haben, daß wir uns seiner Güte und Gnade sollen trösten, wenn wir gleich nicht werth sind noch verdient haben. Er will nicht, daß wir wie Moses und Aaron denken sollen, er werde denn allererst gnädig sein, wenn wir fromm und solcher Hülfe würdig sind. Daß aber Moses ins gelobte Land nicht kömmt, hat seine Deutung; denn Gesetz und Evangelium, Werk und Gnade, sind wider einander; so ist das gelobte Land ein Bild der ewigen Seligkeit, da gehöret nicht Moses, sondern Christus zu, nicht Verdienst, sondern Gnade und Barmherzigkeit.

\section{Votum}

Der barmherzige GOtt mehre unsern Glauben, und verleihe, daß wir den durch gute Werke leuchten lassen, und gebe uns das Ende des Glaubens, das ewige Leben, Amen.
\quellseite{254}

\chapter{Das 21. Kapitel}

\section{Vorrede}

Das Bild der ehernen Schlange, welche Moses nach Gottes Befehl in der Wüste erhöhet hat, soll allen Christenherzen wohl bekannt sein. Derwegen laßt uns dessen Deutung und Erklärung anhören. Dieses Bild aber erkläret der HErr Christus selbst, daß er dadurch für gebildet und gemeinet sei, wie zu lesen Johannes am 3.: Denn also saget er: (Gleich wie Moses in der Wüste eine Schlange erhöhet hat, also muß des Menschen Sohn auch erhöhet werden, auf daß alle, die an ihn glauben, nicht verloren werden, sondern das ewige Leben haben. Allhier ist zu merken, worinnen denn die erhöhete Schlange eine Figur und Deutung ist Jesu Christi, nämlich darinnen: Erstens, gleich wie die eherne Schlange wohl allerdings eine Form und Gestalt hat einer rechten, wahrhaftigen Schlange, aber doch ist sie ohne alles tödtliche Gift gewesen: also hat Jesus Christus wohl aller Dinge menschliche Natur und Affecten an sich, doch ist er ohne Sünde und Gift der Sünden, wie St. Paulus sagt 2 Cor. 5.: Der von keiner Sünde wußte, Jesus Christus, den hat Gott für und zur Sünde gemacht, auf daß wir in ihm würden die Gerechtigkeit, die für Gott gilt. Und Röm. 8.: Gott sandte seinen Sohn in der Gestalt des sündlichen Fleisches, und verdammte die Sünde durch Sünde, auf daß die Gerechtigkeit, im Gesetz erfordert, in uns erfüllet würde. Darnach in dem Stücklein ist auch die erhöhete Schlange ein Vorbild Jesu Christi: daß gleich wie der, der von den rechten feurigen Schlangen gebissen war, und doch diese erhöhete eherne Schlange mit fröhlichem Vertrauen hat können ansehen, der ist von dem tödtlichen Gift entledigt und genesen: also wer diesen am Kreuz und durch seine Auferstehung erhöheten Jesum Christum im Glauben wird ansehen, dem soll der alten Schlangen, das ist des Teufels und Todes giftiger Biß, sogar nicht schaden, daß auch alle, die an ihn glauben, nicht sollen verloren werden, sondern haben das ewige Leben. Doch ist auch heilsam dies zum Unterschied zu wissen: Moses Schlange hat die Kraft zu heilen und Gift zu vertreiben, nicht ab intra, sed ab extra, das ist, nicht in und von sich selbst, sondern von dem lebendigen Wort Gottes, der solche Kraft dadurch hinein gesprochen hatte; Jesus Christus aber hat die lebendige Kraft in und von sich selbst, und ist uns fürgestellet zu einem Gnadenstuhl durch den Glauben in seinem Blut, Röm. 3.

\section[Der Biß feuriger Schlangen, nach Anschauung der ehernen]{Der Biß feuriger Schlangen, nach Anschauung der ehernen Schlange geheilet}

a Harma) heißt ein Bann.

b Feurige) darum heißen sie feurige, daß die Leute von ihnen gebissen, durch ihr Gift feuerroth wurden, und für Hitze starben, wie an der Pestilenz oder Carfunkel.

c Arnon) ist derselbe hohe Fels, der in Besupha, das ist in Wolken, Winden und Wetter stehet. Unten fließen Bäche, und er an denselben seinen Bächen hin währet bis gen Ar.
\quellseite{255}

d Stäbe mögen auch die heißen die Fürsten selber; wie im Zeichen der Könige in Egypten ein Stab heißt, darum, daß sie das Volk regieren, wie ein Hirte das Vieh mit seinem Stabe.

e Töchtern) das ist, die Dörfer und Flecken um die Stadt gelegen.

\section{Summarium}

Ist eine feine Historia, wie das Volk wider Mosen murrete, und verdrossen und unwillig ward, daß sie so lange in der Wüste bleiben mußten, V. 5. Solchen Unwillen strafet Gott mit giftigen Schlangen; aber da Moses für sie bittet, befiehlt Gott, daß man eine eherne Schlange machen soll, und verheißet, wer die ansehe, dem soll das Gift nicht schaden, V. 9.

Solche Historia deutet unser lieber HErr Christus Joh. 3, 13. auf sich und sagt, die eherne Schlange hab ihn bedeutet, denn gleich wie dieselbe einer rechten Schlange gleich gewesen, aber doch nicht allein sein Gift gehabt, sondern auch wider das Gift geholfen hat. Also ist Christus ein Mensch wie andere Menschen, aber er hat keine Sünde, denn er ist Gott, und vom heiligen
\quellseite{256}
Geist empfangen, und von einer Jungfrau geboren. Er hilft aber wider das Gift der Sünden allen denen, so ihn mit Glauben am Holz des Kreuzes ansehen, daß er für ihre Sünden bezahlet und genug gethan habe.

Sie hebt sich auch die Historia an, wie die Kinder von Israel wider die Heiden jenseit des Jordans gestritten, und ihre Länder eingenommen haben.

\section{Votum}

Lob, Ehr, Preis und Dank sei dem HErrn JEsu Christo für alle seine Wohlthaten, die er uns durch seinen Gehorsam und Kreuz erworben hat; der wolle uns an Leib und Seele heilen, und alle selig machen, Amen.

\chapter{Das 22. Kapitel}

\section{Vorrede}

Der Geiz wird von den Alten genennet eine Stiefmutter der Gerechtigkeit, und eine Hauptsach aller Laster. St. Paulus nennet den Geiz eine Abgötterei und eine Wurzel, das ist, einen Ursprung und Quellbrunn alles Uebels, 1 Tim. 6. Die Alten haben denselben Geiz abermals genennet eine unersättliche Trunkenheit, und ein solch Laster, das bei und in den Menschen nimmer alt werde; denn gemeiniglich, je älter einer wird, je karger und geiziger er wird, und je näher die Leute zur Herberge (das ist, zum Grabe) haben, je ernstlicher sie nach der Zehrung pflegen zu trachten. Derwegen sich niemand den Geiz einnehmen soll lassen, sondern ein jedes soll den meiden und meiden, in Betrachtung, daß Gott für uns sorgen, und keinen Mangel uns will leiden lassen. Denn Gottes Ordnung ist für und für, wo er das Leben gibt, daß er auch schaffet, daß man es erhalten könne. Darauf höret auch an das 22. Kapitel.

\section{Bileam soll den Israeliten fluchen. Seine Eselin redet}
\quellseite{257}

\section{Summarium}

Ist eine Historia, wie der Moabiter König den Propheten Bileam holen ließ, daß er den Juden fluchen sollte, V. 5. Sie ist ersichtlich zu merken, daß es scheinet, als sei Bileam ein rechter Prophet gewesen, der nicht mit Zauberei umgangen, sondern aus dem Geist Gottes geweissaget habe; wie aus seiner Weissagung unten wohl gespüret wird, und hie klar stehet, daß Gott mit ihm geredet habe.

Zum andern, daß er Gott zweimal fraget, V. 11. 19., und Gott ihm erstlich heißt, er soll nicht ziehen, V. 12., und zum andern Mal heißt, er soll ziehen, V. 20. Das ist also zu verstehen, daß Bileam am ersten Befehl und an der ersten Antwort sich sollte haben genügen lassen, und weiter nichts gefraget haben. Denn Gott ist in seinem Wort beständig. Darum was er einmal sagt, dabei soll man bleiben. Daß aber Bileam den ersten Befehl verachtet, und zum andernmal fraget, thut er unrecht, und versucht Gott um seines Geitzes willen. Darum versucht ihn Gott wieder, und heißt ihn ziehen. Denn wer sich Gottes Wort nicht will warnen lassen, der mag ihm selbst zum Aergerniß und Schaden fortfahren, und sehen, wie es ihm endlich gerathe. Gott soll man nicht versuchen. Es heißt aber Gott versuchen, wo man bei einem Wort nicht bleibet, und nicht einfältiglich glaubt, sondern weiter fragt, und denkt, Gottes Wort sei wie eines Menschen Wort, das sich ändern und beugen läßt.
\quellseite{258}

\section{Summarium}

Ist eine wunderbärliche Historia. Nachdem der Prophet Bileam sich des ersten Befehls Gottes, um seines Geizes und Menschengunst willen, nicht hat gehalten, sondern Gott versucht, und zum andernmal gefragt, strafet ihn Gott; wie es St. Petrus in seiner andern Epistel bezeuget durch die stumme Eselin, darauf er reitet, daß er unrecht habe gethan, daß er um Geldes und unrechten Lohnes willen sich wider den ersten Befehl Gottes hat lassen aufbringen, gleich wie ihn der Engel auch strafet. Aber wer einmal Gottes Wort nicht hat wollen folgen, der ist wohl werth, daß er ihm selbst zum Schaden und Aergerniß fortfahre, und darüber zu Boden gehe. Wie dieser Ursach halber diese Historia beschrieben ist, uns zur Warnung, ob wir gleich hohe Gaben haben, daß wir doch beim Wort auf das einfältigste bleiben, und uns davon nicht sollen abweisen lassen.

\section{Votum}

Gott regiere uns, daß wir ihn ja nicht versuchen, sondern daß wir seinem Wort gehorsamlich folgen, für Geiz und andern Sünden uns hüten, damit wir von ihm gesegnet werden, zeitlich und ewiglich, an Leib und Seele, Amen.

\chapter{Das 23. Kapitel}

\section{Vorrede}

Wem Gott wohl will, dem können der Teufel und böse Menschen nicht schaden. Denn also sagt Christus, Matth. 16.: Die Pforten der Hölle sollen sie (vernehmt meine Gläubigen) nicht überwältigen. Und David im 118. Psalm spricht: Der HErr ist mit mir, darum fürchte ich mich nicht, was können mir Menschen thun? Derwegen erwarte ein jedes seines Berufs und befehle sich täglich in Gottes Schutz und Segen, und lasse den Teufel Hulfstricke legen, und die bösen Menschen suchen und Unglück stiften, wie sie wollen und mögen; Gott behütet doch die Seinen, und wo die bösen Flüche ausgehen, da geben sie wieder ein, mit Unglück und Schaden. Dies vermahnet auch das 23. Kapitel.

\section{Bileams Fluch in einen Segen verwandelt}

a Kirchen Baal.

b (Herr) hie leugt Bileam, daß er will zum Herrn gehen, welcher ihm bereits gesagt hatte, daß er nicht sollte gehen; sondern er gehet zu seinen Zauberern unter des Herrn Namen.

c (Spruch) heißt die Oracula, das ist, solche Worte, die er nicht von ihm selbst redet; sondern die ihm Gott in den Mund gab. Als wenn ein Gottloser den Text des Wortes Gottes spricht, daß wider ihn selbst und die Seinen ist.
\quellseite{259}

d Menschen lügen, und müssen auch zuweilen fehlen, da sie nicht halten können, denn sie sind selbst ihres Lebens ungewis.

e Mühe und Arbeit heißt die Schrift Abgötterei, oder falschen Gottesdienst, und was ohne Glauben geschieht, Ps. 10, 7. Unter seiner Zunge ist Mühe und Arbeit.

f Trommeten des Königs) das ist, die lieblichen Trommeten Gottes ihres Königs, der sie zu machen befohlen hatte; darum sie unüberwindlich waren im Streit. Mag auch wohl heißen das Wort Gottes, so in diesem Volke lauter und öffentlich gelehret ward.

\section{Summarium}

Dies ist eine feine Historia, in welcher man sieht, wie es allen Menschen und dem Teufel selbst unmöglich ist, wo Gott seinen Segen gibt, daß man denselben könne wenden oder ändern. Wie denn Bileam nicht anders kann, denn daß er die Kinder Israel segne, V. 8. 20. Darum wer in Gottes Furcht und seinem Wort und Willen lebt, der darf sich für nichts besorgen. Denn ob die Leute schon wollten Schaden thun, so ist doch Gott auf jener Seite und behütet.

Sonderlich aber ist hie zu merken, daß Bileam die Juden fürnehmlich darum lobet und segnet, daß weder Mühe noch Arbeit, das ist, keine Abgötterei noch falscher Gottesdienst, sondern Gott und seine Trommeten, das ist, sein Wort unter ihnen ist, V. 21. Wo das ist, da darf man sich für Zauberei und Fluchen, ja für dem Teufel selbst nicht fürchten; wie es die Erfahrung noch heutiges Tages mit bringt. Wo man aber Gott nicht fürchtet, an sein Wort sich nicht hält und in Sünden lebt, da ists nicht Wunder, ob schon aus Gottes Verhängnis der Teufel, oder böse Leute, Schaden anrichten; denn sie können ohne das anders nichts, denn Schaden thun.

\section{Votum}

Der allmächtige Gott steure und wehre dem Teufel und bösen Menschen; erhalte uns in seiner Furcht, und behüte uns für allem Unglück, an Leib und Seele, bis zum ewigen, seligen Leben, Amen.

\chapter{Das 24. Kapitel}

\section{Vorrede}

Unglauben und Geiz fallen noch heute viele Menschen, die da meinen, sie werden selig werden, und werden doch im Abgrund der Hölle in
\quellseite{260}
Einigkeit müssen verlieb nehmen, und Qual leiden. Dies lehret der heilige Täufer Johannes, da er spricht, Joh. 3.: Wer dem Sohn nicht glaubet, der wird das ewige Leben nicht sehen, sondern der Zorn Gottes bleibet über ihm. Und St. Paulus Eph. 5. sagt also: Kein Geiziger, welcher ist ein Götzendiener, hat Erbe an dem Reiche Christi und Gottes. Darnach schaue ein jeder zu, wie er sich in Kaufen und Verkaufen, und in allerlei Handel und Wandel halte, und wie und was er glaube, und niemand betrüge sich selbst. Denn was Gott an andern gestraft hat, das wird er ihm an mir und euch nicht gefallen lassen. Davon lehret auch das 24. Kapitel.

\section{Bileams Weissagung von dem Stern aus Jakob}

a Hieraus merket man, daß Bileam droben allzeit zur Zauberei gegangen, unter Gottes Namen; aber der HErr ist ihm immerdar begegnet und hat die Zauberei gehindert, daß er hat müssen das rechte Wort Gottes fassen und aussprechen, anstatt der Zauberei.

b Amalek war der erste unter den Heiden, so die Kinder Israel angefochten, 2 Mos. 17, 8., aber durch Saul vertilget, 1 Sam. 15, 7.
\quellseite{261}

c Chittim) sind die aus Europa, als der große Alexander und Römer, welche auch zuletzt untergehen. Und zeuget hie die Weissagung, daß alle Königreiche auf Erden eins nach dem andern untergehen müssen, neben dem Volk Israel, welches ewig bleibt um Christi willen.

\section{Summarium}

Dies ist eine Weissagung des Propheten Bileams, wie seine Macht den Kindern Israel soll fürstehen, finstermal Gott bei ihnen ist, und will sie gesegnet haben, V. 4.

Daß er aber meldet von dem Stern aus Jakob und dem Scepter Israel, V. 17., das geht auf den Herrn Christum, welches Reich alle andre Reiche dämpfen, und alles unter sich bringen soll, nicht mit dem Schwerdt oder weltlicher Gewalt, sondern allein durchs Wort.

Aber ehe Christus kommt, sagt er, sollen die Assyrer gewaltig werden, und alles unter sich bringen, bis der große Alexander aus Chittim die Assyrer auch dämpfe, und endlich alle Gewalt an die Römer komme, und Christus geboren werde, V. 24. Solches ist eine treffliche Weissagung, und ohne Zweifel aus dem h. Geist. Darum Bileam ein rechter Prophet gewesen und Gott erkennet hat, aber durch Geiz und Unglauben ist er gefallen, wie andere große Leute mehr. Denn Weissagen ist eine solche Gabe des h. Geistes, die wohl Leuten widerfahren kann, die dennoch nicht selig werden.

\section{Votum}

Der treue, barmherzige GOtt regiere uns alle, daß wir für Sünden uns fleißig hüten, und im wahren Glauben gegen ihn, und in christlicher Liebe gegen unsern Nächsten beständig leben und selig sterben, Amen.

\chapter{Das 25. Kapitel}

\section{Vorrede}

Wie alle Menschen Abgötterei und Unzucht fliehen sollen, da sie nicht Leib und Seele zu büßen wollen: also soll die Obrigkeit Abgötterei abschaffen, und die Unzucht ernstlich strafen, damit Gott Land und Leute derwegen nicht strafen dürfe. O beiderseits eine nöthige Warnung. Demwegen laßt uns nicht leben, daß uns dieß Leben gereue, und uns und andern zum Verderben gelange; sondern laßt uns leben, daß wir um Christi willen alhier und in Ewigkeit in Gottes Gnade und selig leben, und daß wir unsern Nächsten durch Kraft des heiligen Geistes viel nützen und dienen. Darauf höret ferner an das 25. Kapitel.

\section{Abgötterei und Hurerei wird ernstlich gestraft}
\quellseite{262}

\section{Summarium}

Es ist ein schrecklich Exempel, wie das Volk Israel der Heiden Abgötter anbetet, und mit den heidnischen Weibern Hurerei treibet, V. 2. Solche Sünde strafet Gott mit einer so grausamen Plage, daß ihrer vier und zwanzig tausend darüber todt blieben, V. 4.; auf daß wir lernen, wie Gott beiden Sünden, der Abgötterei und Hurerei feind sei, und wolle sie greulich strafen.

Darnach ist ein nützlich Exempel mit Pinehas, des Aarons Enkel, der den Hurer sammt der Hure ersticht, V. 7. Denn da sehen wir, wo die Obrigkeit solche Sünde mit Ernst strafet, daß Gott die gemeine verdiente Strafe abwenden will, welche Strafe sonst müßte folgen, wo dergleichen Sünden von der Obrigkeit nicht gestraft würden. Denn wo die Obrigkeit mit Strafen faul und nachlässig ist, da wird Gott verursacht, die Sünde zu strafen, und gebet gemeiniglich darnach solche Strafe über den ganzen Haufen, wie man in Exempeln siehet.

\section{Votum}

Gott regiere uns alle durch seinen heiligen Geist, daß wir uns für Abgötterei und aller Unzucht hüten, und helfe, daß die Obrigkeit alle Laster ernstlich strafe, ihm zu Ehren, und den Menschen zur Besserung, Amen.

\chapter{Das 26. Kapitel}

\section{Das jüdische Volk wird von neuem gezählet}

a) Zeichen) ein Schreck-Zeichen, daran sie gedächten und sich für gleicher Sünde hüteten.
% No apparatus blocks retained on this page. See selection.json for block-by-block exclusion reasons.
\quellseite{264}

\chapter{Das 27. Kapitel}

\section{Vorrede}

Wir Menschen sind zur Ordnung von Gott erschaffen; aber keine Kreatur oder Geschöpf auf Erden lebet unordentlicher, als die Menschen. Darüber klagt der heilige Prophet Jeremias im 8. Kapitel, und sagt: Ein Storch unter dem Himmel weiß seine Zeit, eine Turteltaube, Kranich, Schwalbe, merken ihre Zeit, wenn sie wiederkommen sollen; aber mein Volk will das Recht des HErrn nicht wissen. Weil aber Gott gewisse Stände geordnet hat, und jedem Stande in seinem Wort Regeln und Weise vorgeschrieben, wie ein jeder Mensch in seinem Stande christlich und unsträflich leben soll; deswegen laßt uns Gottes Ordnung treulich folgen, damit er uns aus Gnaden um Christi willen, nach diesem elenden Leben, in das ewige und freudenreiche Leben ordne und versetze. Darauf höret auch an das 27. Capitel.

\section[Gesetz von Erbtöchtern. Josua an Mosis Statt zum]{Gesetz von Erbtöchtern. Josua an Mosis Statt zum Fürsten des Volks geordnet}
\quellseite{265}

a Deine Herrlichkeit) das wird vielleicht eine sonderbare Weise gewesen sein, daß Mose dem Josua den Stab oder die Hände aufs Haupt gelegt hat. Gleich wie man die Könige salben, oder, wie man die Lehen zu empfahen pflegt; welches alles muß eine Weise und Gepränge haben.

b Das Licht) das ist, das Licht auf der Brust des Hohenpriesters, 2 Mos. 28, 30. Daher sagen etliche, wenn Gott habe auf des Priesters Frage geantwortet, das hat sollen Ja sein, so habe das Licht Glänzen von sich gegeben. Es haben aber hernach die Könige also gebraucht, wenn sie Gott um Rath fragen; als 1 Sam. 28, 6. c. 30, 7.

\section{Summarium}

Am ersten ist eine seine nöthige Ordnung, wie es mit den Testamenten oder mit dem Bescheiden der Güter soll gehalten werden, wenn sich ein Todsfall begibt. Da ordnet es Gott also, daß das Gut allweg auf den nächsten Bluterben fallen soll, V. 6. 11.

Darnach ist ein sehr fein Exempel, als Gott Mosen abfordert von diesem Leben, und sagt, er müsse sterben, V. 13., da ist sein höchstes Anliegen und Begehren, daß Gott seinem Volk einen andern Fürsteher ordnen wolle, damit sie nicht ohne Regiment, wie die Schafe ohne einen Hirten seien, V. 16. Solchem Exempel soll alle Obrigkeit nachfolgen, und für ihre Unterthanen auch also sorgfältig sein und bitten.

\section{Votum}

Helfe der liebe GOtt, daß beide, Obrigkeit und Unterthanen, seinem Wort sich gemäß halten, in christlicher Regierung und gehorsamer Unterthänigkeit, und beschere uns allen zu seiner Zeit ein seliges Sterbestündlein, Amen.

\chapter{Das 28. Capitel}

\section{Gesetz von mehrerlei Opfern wiederholet}
\quellseite{266}

a Wein) ist hie sicera, hoc est, inebriativum, id est, purum, non dilutum, sed merum et efficax. Quia verbum purum inebriat vere animas. Das ist, berauschendes Getränk, d. i. rein, nicht verdünnt und geschwächt, sondern lauter und kräftig. Weil das reine Wort die Seelen wahrhaft trunken macht.

\chapter{Das 29. Capitel}

\section{Noch andere Opfer}
% No apparatus blocks retained on this page. See selection.json for block-by-block exclusion reasons.
\quellseite{268}

a Ezere(th) als auf dem Kirchhof versammelt und beschlossen, ut liceat audire maledictiones etc., daß man die Versuchungen hören könne. 5 Mos. 28, 15.

\chapter{Das 30. Kapitel}

\section{Von Gelübden, wie sie verbinden, oder nicht verbinden}

a Ihre Seele) das ist, wenn sie sich verbünde zu fasten, oder sonst etwas zu thun, mit ihrem Leibe (Gott zu Dienst; daß Seele hie heißt so viel als der lebendige Leib, wie die Schrift allenthalben braucht.

\chapter{Das 31. Kapitel}

\section{Vorrede}

Ach der großen Sicherheit und Sünden dieser letzten Welt. In allen Predigten wird gesagt vom Zorn Gottes wider die Sünde; von dem zeuget Gott selbst auch, erschrick vom Himmel durch Wunderzeichen, und durch ungewöhnliche Wetter, dann auf Erden durch allerlei wunderliche Unfälle. Wer fürchtet sich aber für Gottes Zorn? Wer bessert deswegen sein Leben? Sondern es gehet wie der heilige Prophet Jeremias im 5. Kapitel klaget, und sagt: HErr du schlägest sie, aber sie fühlens nicht, du plagest
\quellseite{269}
ke, aber sie bessern sich nicht, sie haben ein härter Angesicht denn ein Fels, und wollen sich nicht bekehren. Derwegen laßt uns Buße thun, und also leben, daß Gott nicht vollends seinen Zorn wie Wasser auf uns ausgieße, und durch seine Strafen, in zeitlichen und ewigen Verderb stürze. O HErr Christe, das hilf uns allen, Amen. Darauf höret nun an das 31. Kapitel.

\section[Israel erhält wider die Midianiter den Sieg, und]{Israel erhält wider die Midianiter den Sieg, und große Beute}

\section{Summarium}

Hie ist eine Historia, wie das Volk Israel die Midianiter geschlagen, und ihr ganzes Land geplündert haben, darum daß sie das Volk Israel zur Sünde gereizt und verführt hatten. So sieht man den Zorn Gottes wider die Sünde, daß er alles heißt erstechen und erwürgen, ausgenommen die Mägdlein, so noch nicht mannbar waren, der waren an der Zahl zwei und dreißig tausend, v. 7. 35., und zum Exempel und zur Besserung, auf daß wir lernen, daß Gott endlich pflegt die Sünde zu strafen, wo man nicht ablassen noch sich bessern will.
\quellseite{270}

\section{Summarium}

Dies ist ein fein Exempel, wie Gott befiehlt, daß man von der Beute, so an den Midianitern gewonnen war, den Leviten auch ihren Theil geben soll zum Zeugniß, daß der Sieg vom HErrn kommet, und die Leute denselben brauchen sollen zur Ehre Gottes und Förderung des Gottesdienstes. Wie am Ende die Hauptleute auch thun, und eine große Summe Goldes schenken zur Hütte, v. 48., zum Bekenntniß und Dank, daß ihnen Gott den Sieg gegeben und sie für Unglück behütet hätte. Das sind keine gottesfürchtige Kriegsleute gewesen, darum

15*
\quellseite{271}
hat ihnen Gott auch Glück gegeben. An solchen mangelt es uns wider den Türken. Denn unser Kriegs-Volk fragt gemeiniglich nach der Kirche und predigt nicht viel, spielen, Unzucht treiben, plündern und rauben, gottloslästern und fluchen ist ihr Brauch. Darum ist weder Glück noch Heil bei ihnen, Gott erbarme es.

\section{Votum}

Gott erbarme es, und streite selber für seine arme Christenheit wider den Türken, und alle andere Feinde, und erhalte uns den Land- und Stadtfrieden gnädiglich, bis an den lieben jüngsten Tag, Amen.

\chapter{Das 32. Capitel}

\section{Vorrede}

Gott hilft und segnet darum viel Leute, daß die so viel haben, denen rathen und dienen sollen, die nichts oder wenig haben. Aber dies kehren die Weltkinder um, und suchen immer weiter ihren Nutzen und Vortheil, wo auch der Nothdürftige bleibe. Derwegen trachte also ein jeder auf sein Bestes, daß er seines Nächsten nicht vergesse, sondern dem wo und wie er kann, helfe. Denn Nothdürftige verlassen, bringt Gottes Fluch und Strafe; dagegen wohl dem, der sich des Dürftigen annimmt, denn wird der Herr erretten zur bösen Zeit, sagt David im 41. Psalm. Ferner höret an das 32. Capitel.

\section{Der Besitzung des Landes Canaan wird ein Anfang gemacht}
\quellseite{272}

\section{Summarium}

Ist eine historia, wie Ruben, Gad und der halbe Stamm Manasse, das Land jenseit des Jordans einnehmen und besitzen. Doch also, daß sie den andern sollten helfen kriegen, und nicht Ursach geben, daß sie verzagt würden, V. 19. 25. Denn wo sie allein ihren Vortheil gesucht und sich der andern weiter nicht hätten wollen annehmen, sagt Moses, daß es eine große Sünde wäre, welche Gott zu seiner Zeit strafen würde. Und hie gehet an, daß Gott seine Verheißung vollstreckt, und sein Volk in der Heiden Lande einsetzet.

\section{Votum}

Dem treuen barmherzigen GOtt danken wir für alle seine Wohlthaten; der wolle seiner Verheißung nach an Leib und Seele uns segnen, und helfen, daß wir unsern Nächsten allzeit willig und nach höchstem Vermögen dienen, Amen.

\chapter{Das 33. Capitel}

\section{Register der Reisen und Lagerstätten des Volks Israel}
\quellseite{273}
\quellseite{274}

\chapter{Das 34. Capitel}

\section[Grenzen des gelobten Landes. Wie und durch welche]{Grenzen des gelobten Landes. Wie und durch welche es auszutheilen}
\quellseite{275}

\chapter{Das 35. Kapitel}

\section{Vorrede}

Das Leben ist des Menschen edelstes und herrlichstes Kleinod in dieser Welt, wie denn der Satan selber bekennt, Hiob 2. und saget: Alles was der Mensch hat, läßt er für sein Leben. Nun ist aber die Welt nichts anders, denn ein Stall voller Mörder, die andere tödten mit Herzen, Geberden, Worten und Werken. Derwegen verbeut Gott ernstlich allen Todtschlag, und befriedet uns als der Liebhaber aller Menschen wider aller Leute Frevel und Gewalt, und will, daß keiner dem andern an Leib und Leben schaden solle, und da es jemand thäte, dräuet er, daß er wieder solle gestraft werden an Leib und Leben zeitlich, und da er ohne Buße stirbt, auch ewiglich. Dies lehret auch das 35. Kapitel.

\section{Von den Städten der Leviten, Freistädten und Todtschlag}

Messen, das ist Geometrica proportione gerechnet, nämlich, die Vorstadt soll rings umher, an der Stadt, tausend Ellen weit sein, und eine jegliche Seite der Stadt zwei tausend Ellen lang. Das heißt auf deutsch: Die Vorstadt soll halb so weit sein, als eine Seite der Stadt lang ist. Sie sei viereckt und viereckt, oder wie sie kann, so soll man sie messen und in vier Seiten theilen; und darnach sie groß und klein ist, wird die Vorstadt auch groß oder klein.
\quellseite{276}

\section{Summarium}

Das ist eine Ordnung, daß man soll etliche Freistädt(e) machen für sie, so einen Todtschlag begangen haben, V. 6., nicht erlauben, daß man die Todtschläger befreien sollte, sondern daß man ordentlicher und redlicher Weise sich erkunden möchte, ob der Tod fürseßlich oder ohngefähr geschehen wäre, und also unschuldiges Blutvergießen verhütet würde.

Und ist hie sonderlich zu merken, daß Gott sagt, das Land werde durch Todtschlag verunreinigt, und solche Sünde möge nicht versöhnet werden, denn der Thäter werde von der Obrigkeit auch getödtet, V. 34.; denn man werde mit Geld noch andern, solche Sünde soll verrichten, noch versöhnen lassen. Und wo gleich ein Todtschlag nicht fürseßlich, sondern ohngefähr sich begeben hätte, müßte dennoch derselbe Thäter flüchtig sein, und sich in der Freistädte einer enthalten, V. 24., und dürfte nicht ehe zu seinem Erbe kommen, bis der Hohepriester stürbe, V. 24. Also hat Gott mit sonderlichem Ernst über dem fünften Gebot gehalten, und soll weltliche Obrigkeit solchen Ernst wohl bedenken, und derhalben ihrem Amt nach keinen Todtschlag ungestraft noch abbitten lassen; denn unschuldiges Blut schreiet zu Gott, und wo es an dem Rechtschuldigen nicht gerochen wird, muß der ganze Haufen und ein ganzes Land solche Sünde tragen und darum verhalten.

\section{Votum}

Wir danken Gott, dem leutseligen Menschenfreunde, daß er unserm Leib und Leben beschützet; der steure allen mörderischen Menschen, und regiere uns, daß wir niemand an seinem Leib und Leben weder Schaden noch Leid thun, Amen.

\chapter{Das 36. Kapitel}

\section{Vorrede}

Gott hat im jüdischen Volke die Erbe und Güter unterschieden, und geboten, daß die liegenden Gründe von einem Stamm an den andern
\quellseite{277}
 nicht erben sollten; denn Gott wollte, daß jedermann wissen sollte, aus welchem Stamm er wolle, daß der verheißene Messias sollte geboren werden. Daraus folget, daß Gott einem jeden das Seine gönnet, daß er mit gutem Gewissen erworben oder ererbet hat, und daß er nicht wolle, daß alles gemein sein sollte, wie davon die Wiedertäufer lehren. Und derwegen auch ein jeder Christ mit gutem Gewissen sein Eigenes haben und besitzen kann, doch daß er sein Herz auf den ungewissen Reichthum nicht stelle, sondern auf den lebendigen Gott traue, und daß er das, so ihm Gott bescheret, auch Gott zu Ehren, ihm selbst und andern zu Nutz anwende und gebrauche. Darauf höret an das 36. und letzte Kapitel.

\section[Durch ungleiches Heirathen sollen die Erbtheile nicht verrückt]{Durch ungleiches Heirathen sollen die Erbtheile nicht verrückt werden}

\section{Summarium}

Dies ist eine feine bürgerliche Ordnung, daß die Töchter, so eigen Erbe hatten, nicht durften in fremde Stämme heirathen, V. 6., auf daß die Theilung des Landes, wie sie Gott geordnet hatte, unverrückt bliebe.

\part{Das fünfte Buch Mose}

\chapter{Das 1. Kapitel}

\section{Vorrede}

Das fünfte Buch Mose hat nach der griechischen Sprache den Namen, und heißt Deuteronomium, zu deutsch, das wiederholte Gesetz. Darum daß Moses in diesem Buch wiederholet, erstlich die fürnehmsten Historien, und darnach fast das ganze Gesetz, welches er hie besser erkläret, und weitläufiger ausleget. Sonderlich aber, was die erste Tafel betrifft, wie man Gott von Herzen fürchten und ihm vertrauen soll, und sich seines Willens und Worts allein halten. In dem ersten Theil dieses Kapitels wird wiederholet die Historia aus dem andern Buch Mose am 18. Kapitel, in welcher sonderlich zu
\quellseite{278}
merken ist, wie Mose die vermahnet, welchen das Regiment befohlen ist, daß sie keine Person ansehen, den Armen wie den Reichen hören, und sich für Niemand scheuen sollen. Denn richten sei Gottes Amt; darum soll man es treulich und rechtschaffen führen, und auf Gott mehr denn auf die Menschen sehen. Solche Vermahnung sollen alle Regenten ihnen lassen sonderlich befohlen sein. Darauf höret nun an das 1. Capitel.

\section{Gottes Gutthaten. Israels Undankbarkeit}
\quellseite{279}

\section{Summarium}

Ist eine Historia aus dem 13. Kapitel des vierten Buches Mose, wie das Volk GOtt nicht vertrauete, daß er bey ihnen sein und für sie streiten würde. Darum strafet sie Gott, daß sie vierzig Jahre in der Wüste bleiben, und daselbst alle sterben mußten, und keiner in das gelobte Land kommen, auf daß wir lernen, wie man Gott nicht höher unehren, noch ehe zu Zorn bewegen kann, denn so man an ihm und seiner Hülfe und Macht zweifelt. Darum tröstet Moses so treulich, und will, daß man doch hinter sich denke, wie Gott in Egypten und am rothen Meer geholfen, und sein Volk, wie ein Vater seine Söhne, getragen und geleitet habe. Aber solches Predigen war gar vergeblich, gleich wie es auch vergebens war, da er erboth, sie sollten wider die Amoriter nicht streiten, und sie darum geschlagen wurden.

\section{Votum}

Helffe der treue, barmherzige GOtt, daß die Obrigkeit ihr Amt recht und glücklich führe; und stärcke unsern Glauben, daß wir ihm allzeit beständig vertrauen, in Noth und Tod, Amen.
\quellseite{280}

\chapter{Das 2. Kapitel}

\section{Vorrede}

Unser HErr Jesus Christus spricht Johannis am 5. Kapitel: Suchet in der Schrift, denn dieselbige zeuget von mir. Und St. Paulus gebeut Timotheo 1 Tim. 4, 13., er solle anhalten mit Lesen der Schrift, und rühmet Röm. 1., wie das Evangelium sei von Gott in der Schrift verheißen. So weiset uns auch St. Petrus mehr denn einmal zurück in die Schrift; damit Christus und beide Apostel uns je lehren die Schrift des Alten Testaments nicht zu verachten, sondern mit allem Fleiß zu lesen und anzuhören, weil sie selbst das Neue Testament so mächtiglich gründet und bewähren durchs Alte Testament, und sich darauf berufen. So wenig nun des Neuen Testaments Grund und Beweisung zu verachten ist, so theuer ist auch das Alte Testament zu achten. Und was ist das Neue Testament anders, denn eine öffentliche Predigt und Verkündigung von Christo, durch die Sprüche im Alten Testament gesetzt, und durch Christum erfüllet? Darauf höret ferner an das 2. Kapitel.

\section{Vom Sieg wider die Amoriter}
\quellseite{281}

\section{Summarium}

Dies ist eine Historia, in welcher Moses wiederholet, wie des HErrn Wort erfüllet sei, in dem, daß alle die, so an Gottes Hülfe verzagt hatten, in der Wüste die vierzig Jahr sind abgestorben, V. 15. Auch meldet Moses hie etliche Völker, so aus ihrem Lande vertrieben, und andere dagegen eingesetzt sind; auf daß wir lernen, wo dergleichen Strafen ergehen, daß es Gottes Werk sei, der endlich Land und Leute pfleget zu vertilgen, wenn seine Gottesfurcht noch Besserung will folgen.

\section{Summarium}

Ist eine Historia, wie Gott seine Verheißung erfüllet, und der Amoriter König geschlagen, und das Volk Israel sein Land einnehmen läßt. Und ist hie sonderlich zu merken, wie Moses den Amoritern Friede anbeut, daß sie unbeschädigt hätten mögen bleiben, wo sie den Juden den Durchzug erlaubt hätten. Aber die Welt, wenn das Stündlein kommt, muß zu ihrem Verderben selbst Ursach geben. Darum geschichts oft, wie mit den Amoritern hie, daß mancher könnte Friede und Ruhe haben, in Würden und Aufnehmen bleiben, aber durch unnöthiges Gezänk und andern Muthwillen machen sie ihnen und andern Leuten Unruhe, und bringen sich selbst in alles Unglück.

Am Ende ist ein sehr feiner Spruch, da Moses spricht: Es war keine Stadt, die sich für uns schützen konnte, der HErr unser Gott gab uns alles für uns, V. 36. Denn wo Gott seine Hand mit anlegt, da muß es fortgehen. Wiederum, wo er die Hand abzieht, da hilft weder Macht, Geld noch Gestreng. Darum soll man für allen Dingen, sonderlich aber, wo man Krieg führet, darnach trachten, daß man Gott zum Freunde habe, so darf man sich für aller Welt Feindschaft und Macht nicht besorgen.

\section{Votum}

GOtt verleihe, daß wir alle in seiner Furcht leben, steure allen unsern Feinden, und sei und bleibe unser GOtt in Noth und Tod, um JEsu Christi, seines lieben Sohnes, willen, Amen.
\quellseite{282}

\chapter{Das 3. Kapitel}

\section{Vorrede}

Gott ist heut so stark, gewaltig und mächtig, als er von Ewigkeit her gewesen. Daß er aber den Feinden seiner Christenheit so viel gönnet zu seiner Christen Schaden und Verderben, das geschieht nicht derwegen, daß er, Gott der HErr, nicht schützen und helfen könne, sondern unsere Sünden sind deß Ursach, als die Gott strafen soll und muß, und daß man Gott nicht recht bittet um seinen Schutz und Hülfe. Dies lehret Jeremias im 2. Kapitel und spricht: Es ist deiner Bosheit Schuld, daß du so gestraft wirst, und deines Ungehorsams, daß du gestraft wirst. Und St. Jakob in seiner Epistel im 4. Kapitel sagt: Ihr bittet und krieget nicht, darum daß ihr übel bittet. Derwegen daß Gott den Türken und andern unseren Feinden steure, so wollen wir uns bekehren, auch in täglicher Buße leben, und Gott in wahrem Glauben anrufen, daß er uns wider unsere Feinde behüten und erhalten wolle. Dies erinnert auch das 3. Kapitel.

\section{Sieg wider Og, den König zu Basan}
\quellseite{283}

a Haus(e) Kirchen oder Tempel.

\section{Summarium}

Ist eine Historia, in welcher Moses wiederholet, wie Gott den König Og zu Basan in der Kinder Israel Hände gegeben hat. Da stehet man wieder einmal, daß der Sieg Gottes Gabe ist. Am Ende ist ein sonderlich Exempel, wie Moses der treffliche und heilige Mann, Gott darum bittet, er wolle ihn mit ins gelobte Land lassen, v. 23. Aber Gott wird zornig auf ihn, und will ihn nicht erhören, v. 26.; auf daß wir am ersten lernen, wie eine schwere Sünde der Unglaube sei, mit welchem sich Moses bei dem Haderwasser versündiget, und Gott also erzürnet; darnach daß wir lernen unsern Willen in Gottes Willen setzen und geduldig sein, wenn wir um etwas bitten, und dasselbige uns nicht widerfähret; sintemal es Gott gegen Mosen und andere Heiligen auch gethan, und sie in der Anfechtung hat stecken lassen.

\section{Votum}

GOtt stärke unsern Glauben, behüte uns für Ungeduld unter dem Kreuze, und steure und wehre allen Feinden seiner armen Christenheit, Amen.

\chapter{Das 4. Capitel}

\section{Vorrede}

Der königliche Prophet David lehret im 1. Psalm, daß dieß wahre Christen seien, die da Lust haben zum Gesetz des HErrn, und reden von seinem Gesetz Tag und Nacht. Das Gesetz aber nennet David Gottes Wort, das uns Gott (und damit weise und verständig zu machen) offenbaret, und darinnen er uns sein Wesen und gnadenreichen Willen zu erkennen gegeben hat; daraus wir auch lernen, wie wir christlich leben, selig sterben und fröhlich am jüngsten Tage zum ewigen Leben auferstehen sollen. Derwegen soll ein jeder Christ zu diesem Gesetze und Wort Gottes Lust, Freude, Gefallen allezeit und sein Leben lang haben, damit für und für umgehen, und sich darinnen erlustigen, auch solches zu Besserung seines Lebens und Stärkung seines Glaubens fleißig gebrauchen. Ach das helfe uns allen die heilige hochgelobte Dreifaltigkeit, Amen. Darauf höret ferner an das 4. Capitel.

\section{Vermahnung zum Gehorsam des Gesetzes}

a Bewahren) denn Menschen-Lehre hindert Gottes Gebot und führet von der Wahrheit, Tit. 1, 14.
\quellseite{284}

\section{Summarium}

Bisher hat Moses etliche Geschichten wiederholet. Jetzt fähet er an, wie ein treuer Prediger, weil er doch sterben und sein Volk hinter ihm lassen muß, prediget und vermahnet, und warnet sein Volk, daß sie erstlich bei Gottes Wort fest bleiben, darnach demselben ernstlich folgen, und endlich der gnädigen Hülfe und wunderbarlichen Werke Gottes nicht vergessen, sondern dieselbigen ihren Nachkommen predigen und kund thun sollen, V. 2.

Sonderlich aber meldet er, wie Gott am Berge Sinai sich offenbaret und selbst seinen Bund und Gesetz dem Volk fürgetragen habe; demselbigen nach sollen sie sich halten, und keine anderen Götter anbeten, auch kein Bildniß einiger Creatur machen, dieselbigen anbeten, und ihnen zu dienen, wie solches im ersten Gebot ausdrücklich befohlen ist, V. 10.

Wir sollen aber merken, daß Moses die Bildnisse ferner nicht verbeut, denn daß man sie nicht soll anbeten. Sonst kann man ohne Sünde Bildnisse haben zur Erinnerung, wie die Juden die eherne Schlange, Aarons Ruthen, die Cherubin und anderes hatten, und wir in unsern Kirchen, das Bildniß des Leidens, der Auferstehung Christi und anderer Historien haben, nicht zum Anbeten, sondern zur Erinnerung. Solches ist recht gut und nützlich, und in der
\quellseite{285}
heiligen Schrift nirgend verbieten, wie doch die Bildstürmer ohne Grund fürgeben, und ohne allen Unterschied allerlei Bildniß aus der Kirche werfen, so doch wahrhaftige Bilder gleich als eine Schrift sind, welche die Ungelehrten lesen, und sich davon bessern können.

b Das ist, um euretwillen.

c Faciebus suis, scilicet, nube et columna. Das ist, nämlich mit der Wolke und Feuersäule.
\quellseite{286}

\section{Summarium}

Ist eine ernste Warnung, weil Gott durch sein Wort ihnen offenbaret, was sie thun und nicht thun sollen, daß sie sich demselbigen nach mit Fleiß sollen halten, oder Gott werde sie um ihrer Sünde willen aus dem Lande vertreiben, und den Heiden dienen lassen. Doch tröstet er, wo man sich bekehre und Buße thue, daß Gott gnädig und sich ihrer wieder annehmen werde. Jerobalben ist solches wie eine Prophezeyung; denn mit den Juden ist es hernach also gegangen: da Israel Abgötterei trieb, kam der König von Assyrien, und führete sie hinweg, und verwüstete das Land. Das Königreich Juda besserte sich ob solcher Strafe nicht, da that der König von Babel ihnen dergleichen. Solches sind ehe Warnungen, daß wir Gottes Wort für Augen haben und dawider nicht thun sollen.

Sonderlich aber rühmet Moses dieses hie für die größte Wohlthat, die sie erkennen und dankbar dagegen sein sollen, daß Gott sie zum eigenen Volk ihm erwählet und ihnen sein Wort gegeben, und sie so wunderbarlich geführet und erhalten habe, auf daß sie sollen Gott gehorsam sein; so werde es ihnen und ihren Kindern wohl gehen, und Glück und Heil haben, V. 29. Solches alles sollen wir zur Lehre annehmen, und uns auch daraus lernen bessern.

\section{Votum}

Der treue, barmherzige GOtt regiere uns durch sein Wort und den heiligen Geist, daß wir alle in seiner Furcht leben, erhalte uns auch das Evangelium Christi, unsers HErrn, und gebe uns ferner seinen Segen, zu zeitlicher und ewiger Wohlfahrt, Amen.

\chapter{Das 5. Capitel}

\section{Vorrede}

Die zehn Gebote Gottes, so zu reden, sind anders nichts, denn ein göttlich Land- und Stadtrecht, darnach sich alle Christen richten sollen. Denn nachdem ihm Gott der HErr außer auf Erden auch eine Stadt und Bürgerschaft, das ist, eine Christenheit sammeln und erhalten will, so will er dieselbige auch mit gewissen Statuten und Satzungen versehen, darnach sich die Bürger in der Stadt Gottes, das ist die Christenheit, richten sollen. Dies Land- und Stadtrecht, welches sind die heiligen zehn Gebote, hat Gott selbst durch seinen göttlichen Finger in zwo steinerne Tafeln geschrieben. Die erste Tafel begreift die ersten drei Gebote und lehret, wie wir uns gegen dem Landfürsten selbst, das ist, gegen Gott dem HErrn halten sollen, mit dem Herzen nach dem ersten Gebot, mit dem Munde oder Worten nach dem andern Gebot, mit dem Werke gegen sein Wort, nach dem dritten Gebot. Die andere Tafel begreift die letzten sieben Gebote, und lehret, was wir unserm Nächsten zu Ehren und Gehorsam, am Leibe, Weibe, Gut, Ehr, Haus und Hausgesinde, schuldig sind. Dieses sollen wir betrachten bei dem 5. Kapitel.

\section{Wiederholung der zehn Gebote Gottes}
\quellseite{287}

\section{Summarium}

Hie fähet Moses an, und wiederholet die Predigt der zehn Gebote, in welcher Gott seinem Volk anzeigt, was man thun und lassen soll gegen jedermann, wenn man recht thun und selig leben will. Da gehöret nun zu, daß man erstlich solche Gebote recht lerne und wisse, darnach daß man Gott um seine Gnade und heiligen Geist bitte, daß er ein solches Herz geben wolle, auf daß wir solchen Willen also mögen vollbringen, und der Sünde widerstehen; denn wo Gottes Geist nicht ist, da wird solcher Wille Gottes nimmermehr recht erfüllet. Wiederum aber haben wir die Verheißung, daß Gott seinen
\quellseite{288}
Geist denen will geben, so ihn darum bitten, Luc. 11.

\section{Votum}

Diese Verheißung erfülle an uns allen der treue, barmherzige GOtt, und helfe, daß wir nach allen seinen Geboten in rechtem Gehorsam allezeit wandeln bis zum ewigen, himmlischen Leben, Amen.

\chapter{Das 6. Kapitel}

\section{Vorrede}

Wir Christen sind gefreiet vom Gesetz Gottes, was anlanget die Verfluchung und Verdammniß; aber mit nichten sind wir frei, was den Gehorsam belangt. Denn das ist der ewige unveränderbare Wille Gottes, daß die vernünftigen Creaturen, Engel und Menschen, dem Gesetz Gehorsam leisten sollen. Und solchen Gehorsam will ihm Gott gefallen lassen, um seines Sohns unsers Herrn Jesu Christi willen, und will alle die, so sich üben in Gehorsam nach dem Gesetz, segnen an Leib und Seele, zeitlich und ewiglich. So betet, o ihr Christenmenschen täglich und herzlich, Gott wolle euch durch den heiligen Geist regieren und leiten, daß ihr in treuem Gehorsam der heiligen zehn Gebote, gegen Gott im Himmel, und gegen den Menschen auf Erden, leben und sterben möget. Wie ihr denn solches auch in der Taufe eurem Gott und Herrn habt versprochen und zugesaget. Dies erinnert auch das 6. Kapitel.

\section{Erklärung des ersten Gebots, von der Liebe GOttes}

*1 Mos. 18, 19. Schärfen immer treiben und üben, daß sie nicht verrosten noch verdunkeln, sondern stets im Gedächtniß und Wort als neu und helle bleiben. Denn je mehr man Gottes Wort handelt, je heller und neuer es wird, und heißt billig je länger je lieber. Wo mans aber nicht treibet, so wird es bald vergessen und unkräftig.
\quellseite{289}

\section{Summarium}

Dies ist eine Vermahnung, da Moses insgemein die zehn Gebote mit beschließet, daß man Gott fürchten, und diese seine Gebote also hören und behalten soll, daß man sie thue und ins Werk bringe, V. 2. Sonderlich aber vermahnet er, wenn sie ins gelobte Land kommen, und da allerlei Segen und Reichthum finden werden, daß sie des HErrn Werk und Gebot nicht vergessen, sondern dieselbigen ihren Kindern schärfen, das ist, stets davon reden, und im frischen Gedächtniß behalten sollen, V. 6., auf daß sie Gottes Huld behalten, und sich nicht versündigen, wie es doch gemeiniglich gehet; denn wenn es mit uns am besten gehet, und wir alles genug haben, geschieht es gemeiniglich, wo wir nicht sonderlich uns in Gottes Furcht halten, daß wir sicher, stolz und muthwillig werden, und Gottes vergessen. Darum will es vonnöthen sein, daß man in Gottes Furcht wandele, und stets sein Wort übe und treibe, wie Moses hie befiehlet.

\section{Votum}

Helffe der treue, barmherzige GOtt, daß wir nach seinen heiligen Geboten gehorsamlich folgen, auf daß uns an Leib und Seele sein väterlicher Segen folge, in diesem gegenwärtigen elenden und künftigen seligen Leben, um Jesum Christi seines lieben Sohnes willen, Amen.

\chapter{Das 7. Kapitel}

\section{Vorrede}

Gott hat Gefallen an allen denen, die ihn fürchten, und solche haben keinen Mangel, und es wird ihnen wohl gehen in der letzten Noth. Aber dagegen wird Gott ein schneller Zeuge sein wider die, so ihn nicht fürchten, und denen wird es endlich übel gehen. Beides lehret Sirach im 1. Kapitel und saget: Wer den HErrn fürchtet, dem wirds wohl gehen, und wenn er Trost bedarf, wird er gesegnet. Wer aber ohne Furcht fährt, der gefällt Gott nicht, und seine Freiheit wird ihn stürzen. Demwegen wollen wir durch Gottes väterliche Hülfe also leben, damit wir nicht Gott und seine Gebote verachten, sondern wollen den HErrn fürchten, und ihm allezeit dienen, treulich und rechtschaffen, Josua 24. Darauf höret auch an das 7. Kapitel.

\section{Freundschaft mit den Heiden und deren Götzendienst verboten}
\quellseite{290}

\section{Summarium}

Sie ist ein sonderlicher nöthiger Befehl, daß das Volk Israel ganz und gar keine Gemeinschaft mit den Heiden habe, v. 2., sondern sie rein ausrotten und verwüsten soll, mit allen ihren Kirchen, Götzen und Bildern, v. 5., um der Ursache willen, daß sie an ihnen auch nicht lernen und gewöhnen Abgötterei zu treiben. Denn es ist gar leichtlich versehen, daß der Teufel uns Menschen in seine Stricke bringet, wenn wir bei und mit den Gottlosen stets umgehen.

Solche Lehre und Vermahnung ist hoch vonnöthen heutiges Tages, auf daß man des Papsts Geschmeiß mit unnöthigen Kirchen und Klöstern, mit unnützen Büchern und lügenhaftigen ärgerlichen Bildern, und was dergleichen mehr ist, aus dem Weg räume, und nicht also \unsicher{bege…} um des künftigen Aergernisses willen, daß folgen kann, wenn man solches Rattengeniße duldet und nicht ausfeget.

\section{Summarium}

Ist eine sehr schöne und tröstliche Predigt, wie man es so reichlich werde genießen, wenn man Gott fürchtet und seinen Geboten und Worten folget, nämlich, daß Gott Segen des Leibes, beides an Menschen und Vieh geben, alles auf dem Felde reichlich wachsen lassen, allen Krankheiten wehren, Glück und Sieg wider die Feinde geben wolle.

Die sollen wir mit Fleiß lernen und merken, woher doch komme allerlei solcher Unrath, Hunger, Krankheit, Kriege und dergleichen, nämlich davon, daß man Gott nicht fürchtet, und seinem
\quellseite{291}
Wort und Willen nicht folget. Darum sollen wir gern fromm sein, und für dem Bösen uns hüten, weil es Gott so reichlich belohnet, und wo mans nicht thut, so mancherlei Weise strafen will.

\section{Votum}

Unser HErr JEsus Christus, welcher der Weg und die Wahrheit ist, regiere uns, daß wir aller Gottlosen Gemeinschaft fliehen, und helfe, daß wir seinen Geboten folgen, auf daß er auch heut Menschen und Vieh segne und behüte zu unserer Wohlfahrt, Amen.

\chapter{Das 8. Kapitel}

\section{Vorrede}

Wenn es den Menschen übel gehet, so gedenken sie an Gott und beten; aber wenn es ihnen wohl gehet, so vergessen sie gemeiniglich Gottes, und sind undankbar. Solches sollen wir bedenken, und Gott täglich für seine Wohlthaten danken, mit Mund, Herz und christlichem Leben, damit uns der gerechte Gott mit neuem Kreuz und Strafen nicht dazu treiben dürfe. Denn es soll also heißen, wie Gott spricht im 50. Psalm: Rufe mich an in der Noth, so will ich dich erretten, so sollst du mich preisen. Das ist so viel gesagt: Ich stecke dich Mensch darum in Noth, auf daß ich dir Ursache gebe zum Gebet, und ich Ursache habe dir zu helfen, und dir damit auch ferner Ursache gebe, mich zu loben und zu preisen. Dieses soll auch erinnern das 8. Kapitel.

\section[Das Volk wird ermahnet, der Gutthaten GOttes nicht]{Das Volk wird ermahnet, der Gutthaten GOttes nicht zu vergessen}
\quellseite{292}

\section{Summarium}

Ist eine schöne und sehr nöthige Vermahnung, daß das Volk Gottes Gebote halten, und ja seiner Werke und Wunderthaten dann nicht vergessen wolle, wenn sie ins Land kommen, und es ihnen wohlgehet, V. 2. Denn dies ist eine sehr gemeine und gefährliche Anfechtung, daß man unsers HErrn Gottes vergisset und hoffärtig wird, wenn es wohlgehet, darum schickt Gott den Seinen hie auf Erden so mancherlei Trübsal und Anfechtung zu, auf daß er sie in der Furcht und Demuth erhalte, wie Moses hie sagt, V. 16.: Gott habe sein Volk in der Wüste auf mancherlei Weise versucht, auf daß sie seiner nicht vergessen sollten, noch gedenken, sie hätten alles von ihnen selbst.

\section{Votum}

GOtt erbarme sich unser, und gedenke in der Strafe an seine Barmherzigkeit. Er wolle auch verleihen, daß wir an seine Wohlthaten allzeit gedenken, und die loben, preisen und recht gebrauchen, Amen.

\chapter{Das 9. Kapitel}

\section{Vorrede}

Sollte uns Gott nicht mehr zu unserer Leibes- und der Seelen-Nothdurft geben, als daß wir verdienen, ach weh uns, da würde uns nichts mangeln denn alles, und solches überall. Derwegen, weil uns Gott unser Leben lang so große und unzählige Wohlthaten erzeiget hat, und weil er uns noch täglich alle Augenblicke an Leib und Seele Gutes thut, da wir doch eitel Strafe verdienen, so laßt uns Gottes Güte und Barmherzigkeit erkennen, und ihm dafür danken. Dazu vermanet der heilige Prophet Jeremias im 33. Kapitel und saget: Danket dem HErrn unserm Gott Zebaoth, daß er so gnädig ist, und thut immerdar Gutes. Dies lehret auch das 9. Kapitel.

\section{Ursachen der Einnehmung des Landes Canaan}
\quellseite{293}

\section{Summarium}

Dies ist die Vermahnung, wenn sie ins Land kommen und die Heiden vertrieben haben, daß sie Gott die Ehre geben sollen, daß er solches ihnen um seiner Zusagung willen und nicht darum, daß sie es verdienet hätten oder so fromm wären, habe widerfahren lassen, V. 4. Und hie wiederholet er, wie mancherley ihre Väter sich in der Wüste an Gott versündiget haben, auf daß wir lernen, wie wir nicht allein aus Gnaden selig werden, sondern auch alles, was uns Gott in diesem zeitlichen Leben gibt, daß es sei uns aus Gnaden geschenkt, und sich niemand einiges Verdienstes für Gott rühmen dürfe.

\section{Votum}

Lob, Ehr, Preiß und Dank sei GOtt durch JEsum Christum, unsern lieben HErrn, für alle seine großen und vielfältigen Wohlthaten; der wolle uns, die Unseren und die ganze Christenheit, ferner geistlich und leiblich segnen und behüten, als ein treuer, lieber Vater, Amen.

\chapter{Das 10. Kapitel}

\section{Vorrede}

Gott vertrauen und seinem Wort folgen, das sind die Dienste, die Gott von uns fordert, zur Dankbarkeit für seine Wohlthaten, die er
\quellseite{294}
uns reichlich erzeiget. Demwegen weil mir alle bekennen müssen, daß uns Gott ja große und unzählige Wohlthaten bisher mitgetheilet hat, so wollen wir seinem Wort so viel als möglich gehorsam sein, und wollen ihm treulich dienen im Glauben, und unserm Nächsten in wahrer Liebe, daß wir also allzeit dankbar erfunden werden, und daß wir auch mit dem königlichen Propheten David sagen können, aus dem 79. Psalm: Wir dein Volk und Schafe deiner Weide, danken dir ewiglich, und verkündigen deinen Ruhm für und für. So gehöret auch hieher, daß der heilige Apostel Johannes sagt, 1 Epist. Kap. 4.: Das Gebot haben wir von Gott, daß, wer Gott liebet (vernehmt, und die Liebe durch Dankbarkeit beweisen will), daß er auch seinen Bruder liebe. Dazu vermahnet auch das 10. Kapitel.

\section[Die Erneuerung des Bundes soll Israel zur Liebe]{Die Erneuerung des Bundes soll Israel zur Liebe GOttes reizen}
\quellseite{295}

\section{Summarium}

Ist eine Predigt, in welcher Moses etliche mehr Historien wiederholet, und endlich vermahnet, weil Gott sich so gnädig erzeiget, und ihnen nun das Land auch zu eigen geben wolle, daß sie ja wiederum Gott wollen dienen, der anders nichts fodert, denn daß man ihn fürchten und vertrauen, und seinem Worte folgen wolle, V. 12. Insbesondere befiehlt er ihnen die Fremdlinge, daß sie ihnen dienen, und alles Guts thun sollen, weil sie in Egypten auch sind Fremdlinge gewesen, V. 18.

\section{Votum}

Verleihe GOtt, unser lieber Vater, daß sein Wille unser Wille sei, auf daß wir ihm dankbar sein für alle seine Wohlthaten, und daß wir in Gehorsam leben nach seinen Geboten, auch zu Trost unsern nothleidenden Mitchristen, Amen.

\chapter{Das 11. Capitel}

\section{Vorrede}

Was hilfts, wenn die Lehre ist göttlich, und die Menschen leben darauf teuflisch? Also was hilfts, wenn man in der Kirche mit reiner Lehre bauet, und die Leute reißens wieder ein daheim mit gottlosem Leben? Derwegen laßt uns zuschauen, wie wir Gottes Wort hören, und laßt uns dies gebrauchen, auf daß wir sammt den Unsern nicht über unsere Vergessenheit klagen, und über unsern Ungehorsam ewig wehklagen dürfen. Denn wer Gottes Wort verwirft mit Ungehorsam, den wird Gott auch verwerfen mit Zorn, aus seiner Gnade in Noth und Herzeleid; dagegen selig sind zeitlich und ewig, alle die da hören die Worte der Weissagung, und behalten, was darinnen geschrieben ist, Offenb. Joh. 1., und darnach christlich und ehrbar leben. Darauf höret an das 11. Kapitel.

\section[Ursachen, warum GOtt zu lieben und seine Gebote]{Ursachen, warum GOtt zu lieben und seine Gebote zu halten}

a In pedibus eorum: Ich habe gelaufen und gerannt, daß ichs zusammen brächte. Non sterti, nec otiosus fui: Ist mir sauer worden, erjagt, erraunt, erlauffen.
\quellseite{296}

\section{Summarium}

Ist eine sehr schöne Vermahnung, in welcher Moses fast die fürnehmsten Historien wiederholet, in welcher wir sehen, wie GOtt seinen Zorn wider die Gottlosen, und seine Gnade seinem Volk erzeiget hat, V. 3., auf daß sie solche Exempel nicht vergessen, sondern daran lernen sollen GOtt fürchten, und ihm vertrauen, und seinem Wort folgen. Thun sie das (spricht er) V. 13., so werde Gott seinen Segen ihnen reichlich widerfahren lassen; thun sie es aber nicht, so werde sie Gott mit Hunger und Kummer und allerlei Unglück strafen. Solche Vermahnung sollten wir zu Herzen nehmen und uns bessern, so würde Gott nicht also mit unserer Zeit, Krieg, Pestilenz und anderm Unglück unsern Ungehorsam heimsuchen, wie jetzt geschieht und noch mehr geschehen wird, wo wir also in allen Sünden und Ungehorsam fortfahren.

\section{Votum}

Der treue, barmherzige GOtt regiere uns, daß wir von Sünden ablassen, damit wir zu strafen nicht Ursache geben; und helfe, daß wir also leben, auf daß es ihm ein Gefallen sei, und daß er uns aus Gnaden segne, zeitlich und ewiglich, Amen.

\chapter{Das 12. Kapitel}

\section{Vorrede}

Wir danken unserm HErrn Christo herzlich, daß er uns Heiden in seinen Schaffstall und christliche Kirche auch berufen hat, und daß er unser Licht, Trost und Heil sein will. Denn also sagt von Christi Wohlthaten der Prophet Jesaias im 40. Kapitel: Er wird seine Heerde weiden, wie ein Hirte, er wird die Lämmer in seine Arme sammeln und in seinem Busen tragen, und die Schafmütter führen. Und der HErr Christus lehret von ihm selbst auch also, Joh. 10.: Ich gebe meinen Schafen (spricht er) das ewige Leben, und sie werden nimmermehr umkommen, und niemand wird sie mir aus meiner Hand reißen. Lasset uns aber zuschauen, daß wir als seine Schafe seine Stimme und Wort hören, dem glauben, und mit einem
\quellseite{297}
unsträflichen Leben und Wandel folgen, bis zum seligen Todesstündlein. Wie abermals der HErr Christus lehret, Joh. 10. und sagt: Meine Schafe hören meine Stimme, und ich kenne sie, und sie folgen mir. Das helfe uns allen der hochgelobte Sohn Gottes, Christus Jesus, Amen. Darauf höret weiter an das 12. Kapitel.

\section{Ort und Weise des wahren Gottesdienstes}
\quellseite{298}

\section{Summarium}

Hie ist ein ernstlich Gebot, daß die Juden, wenn sie ins Land kommen, ganz und gar nichts sollen lassen bleiben von allem, das die Heiden zu ihrer Abgötterei gebraucht haben, V. 2., allein darum, daß sie nicht Ursach haben gleiche Abgötterei zu treiben. Solches sollen wir auch bedenken, und nichts dulden, was noch vom Papstthum fürhanden ist, und uns zum Aergerniß kann gereichen.

Darnach befehlet Moses insonderheit, wer Gott wolle mit Opfern und andern dienen, daß ers ja nirgend thun soll, denn an dem Ort, welchen ihm Gott selbst dazu erwählen würde, V. 5. Denn Gott kann man nicht dienen, man thue denn den Befehl seines Worts, wie ers ordnet und fordert.

\section{Votum}

GOtt der heilige Geist erleuchte und regiere uns, daß GOttes Wort unserer Füße Leuchte sei, auf daß wir dem in Thun und Lassen gehorsamlich folgen, bis wir durch das finstere Thal des Todes eingehen in den Freudensaal des ewigen Lebens, Amen.

\chapter{Das 13. Kapitel}

\section{Vorrede}

Nachdem Mose mit Fleiß, wie ein treuer Prediger die Seinen vermahnet hat, wie sie Gott fürchten, und in seinem Wort und Willen wandeln sollen, befiehlt er hie, wie man mit den Predigern handeln und umgehen soll, so sie etwas anders lehren und predigen: nämlich, daß man sich an kein Zeichen noch Wunderwerk soll kehren, sondern allein an GOttes Wort halten, V. 3. Was wider das Wort ist, es seien Wunderwerke wie es wollen, soll mans für Teufels-Gespenst und lauter Verführung halten, und solche verführerische Prediger ohne alle Erbarmung tödten, V. 5.

Wo man auf solchen Befehl im Papstthum gesehen, und demselben gefolget hätte, würde der Teufel mit falscher Lehre und falschen Zeichen nicht so viel ausgerichtet, noch die Abgötterei so tief eingerissen haben. Aber weil man die Lehre und das Wort hat lassen fallen, hat der Teufel gute Sache gehabt, und durch lügenhaftige Zeichen den Irrthum wohl gefördert.

\section{Strafe der falschen Propheten, der Verführer und Verführten}

a. Andern) dieser Prophet lehret wider die alte und bestätigte Lehre, und will (wie er sagt, andere Götter) die ersten nicht lassen bleiben, welchem man nicht soll glauben, wenn er gleich Zeichen thät. Aber im Kap. 18, 22. wird der Prophet verdammt, der nicht wider die erste, sondern neue Lehre fürgibt, der soll Zeichen thun, oder nicht gehöret werden.
\quellseite{299}

b Erbarmen, das heißt, Gott über alles lieben.

\chapter{Das 14. Kapitel}

\section{Vorrede}

Gleichwie das vierte Gebot der Ehrenkranz ist, welchen Gott der HErr den Eltern aufgesetzt, gleichwie Gottes Wort und Ordnung des weltlichen Regiments das güldene Band ist, welches Gott der HErr der Obrigkeit hat angehängt, also ist das auch das schöne Kleinod, die güldene Kette und der Freudenbrief, welchen Christus Jesus allen treuen Lehrern und Predigern verehret hat: Wer euch höret, der höret mich, wer euch verachtet, der verachtet mich, Luc. 10. Es sollen aber die Zuhörer ihre christliche Seelsorger, beide ihres Amts und ihrer Person halben, also hören und ehren, daß sie auch dieselben ernähren helfen sollen. Hieher gehöret St. Pauli Vermahnung, die aus dem 6. Kapitel an die Galater also lautet: Der unterrichtet wird mit dem Worte, der theile mit allerlei Gutes dem, der ihn unterrichtet. Irret euch nicht, Gott läßt sich nicht spotten. Das erinnert auch zum Theil das 14. Kapitel.

\section{Trauer-, Speis- und Zehntordnung}
\quellseite{300}

\section{Summarium}

In diesem Kapitel wiederholet Moses etliche Gesetze, so den Juden insonderheit aufgelegt waren. Sonderlich aber treibt er das Gebot mit dem Zehnten, V. 22. und andern, so zu Unterhaltung der Priester und Leviten von GOtt geordnet war, daß man es treulich ausrichten und gern geben soll, wie er droben auch sagt: Hüte dich und verlaß den Leviten nicht, so lange du auf Erden lebest; auf daß wir lernen, wie es Gottes Ordnung und Wille sei, daß jedermann mit dem Seinen gern dazu helfen soll, auf daß die Kirchendiener mögen erhalten werden.

\section{Votum}

Der treue GOtt erhalte uns bei seinem heiligen Worte, und behüte unsere Kirchen und Schulen für falscher Lehre und falschen Gottesdiensten, um der Ehre seines heiligen Namens, und wegen unserer Seelen Seligkeit, Amen.
\quellseite{301}

\chapter{Das 15. Kapitel}

\section{Vorrede}

Freilich, was man den Armen gibt, deß wird niemand ärmer; denn solches wird Gott auf Wucher geliehen, sagt Salomon in seinen Sprüchen im 19. Kapitel. Wohl nun denen an Leib, Seel und Nahrung, ja auch wohl ihren Kindern, ganzem Geschlecht und allen Nachkommen, an geistlichen und leiblichen Gütern, die sich der Armen erbarmen. Denn wahrhaftig des armen Hand ist unsers Herrn Christi Schatzkasten, und ihre eigene Sparbüchse. Derwegen, wer der Armut dienen und helfen kann, der lasse es an ihm nicht mangeln; denn Gott wird ihm wieder dienen und helfen, wo er seiner am meisten bedürfen wird, und wird ihn erquicken auf seinem Sterbebette, Psalm 41. Davon höret auch an das 15. Kapitel.

\section{Vom Erlaßjahr, verkauften Knechten, und Erstgeburt}
\quellseite{302}

\section{Summarium}

Dieß ist eine Ordnung, daß die Juden im siebenten Jahr allwege alle Schuldregister sollen abthun, und nichts mehr an ihre Brüder fordern, V. 2. Weil aber zu besorgen war, sie würden desto weniger mitleiden, wenn das siebente Jahr nahe wäre, vermähnet Mose hie auf das allerschärfste, daß man darum die Armen nicht wolle lassen Noth leiden, V. 8.; denn damit versündiget man sich. Wo man aber gern und willig helfe, da wolle Gott mit reichem Segen kommen, und Glück geben in allem, was man fürnimmt.

Da sollen wir nun auch lernen den Armen die Hand aufzuthun, und gerne helfen, weil es Gott so wohlgefällt, und er so reichen Segen dagegen verheißt.

\section{Votum}

GOtt regiere uns alle, daß wir uns der Armen Noth treulich annehmen, und die es gethan und noch thun werden, die segne GOtt an Leib, Seele und Nahrung, und tröste sie in der letzten Noth, Amen.

\chapter{Das 16. Kapitel}

\section{Vorrede}

Der Mensch ist erschaffen zur Freude in Gott, ist auch nach dem Fall, durch welchen wir in Traurigkeit, Angst und Noth gerathen, wieder erlöst durch den Sohn Gottes, daß er allweg soll fröhlich sein in dem HErrn. So wird der heilige Geist auch gegeben, daß er der Tröster sei der Gläubigen, und in ihren Herzen Freude wirke in betrübtem Zustande. Derjenige welcher Mensch in Noth stecket, oder heimliches Leiden und ein betrübt Herz hat, der ängste sich selber nicht, sondern bete: Gott der Vater wolle ihm sehen lassen sein treues Vaterherz; der HErr Christus wolle ihm zeigen seine heilige tröstliche Wunden; der heilige Geist wolle sein Herz salben mit Freuden-Oel; und sei gewiß, daß ihn Gott segnen, behüten, erretten und erhalten wolle, mit Leib und Seele zur ewigen himmlischen Wonne, Freude und Seligkeit, Amen. Dies erinnert auch das 16. Kapitel.

\section{Von den drei jährlichen Hauptfesten}
\quellseite{303}

\section{Summarium}

Ist eine Ordnung mit den fürnehmsten Festen, Ostern, V. 1., Pfingsten, V. 8. und Laubhütten, V. 13., auf welche alle Mannes-Namen mußten an den Ort kommen, da der Tabernakel oder Tempel war, zum Gedächtnis der Wunderwerke und Wohltat Gottes. Und ist sonderlich zu merken, daß er befiehlet, man soll fröhlich sein für dem HErrn, sintemal er sie so gnädig regieret und geführet, und darnach mit dem Segen so reichlich begnadet hat, V. 11. Solches gnädigen Gottes soll man sich freuen, und allen Trost an ihm haben, und andere Unfälle sich nicht lassen kümmern noch traurig machen.

\section{Votum}

Der treue GOtt tröste und erhalte in wahrem Glauben alle betrübte Herzen in der ganzen Welt, und verleihe, daß wir uns Gottes und seiner Güte und Barmherzigkeit freuen, allweg in Noth und Tod, Amen.

\chapter{Das 17. Capitel}

\section{Vorrede}

Wir danken dem gütigen barmherzigen Gott, der zur Erhaltung der menschlichen Gesellschaft, auch gemeines Nutzes und des Landes und Stadtfriedens, die Obrigkeit geordnet, und bisher zum Zeugniß seiner göttlichen Providenz und väterlichen Fürsorge für uns, erhalten hat. Weil aber durch den Propheten Daniel geweissaget ist, daß am Ende der Welt die Regimente auf schwachen Beinen stehen werden, so bitten wir auch von Herzen, Gott wolle nicht mit uns handeln nach unsern Sünden, und uns nicht vergelten nach unserer Missethat, sondern wolle die wohlverdiente Strafe lindern, und wolle christliche Regimente und Regenten durch seine Gnade geben und erhalten, damit ihm eine Kirche gesammlet, und gemeiner Frieden erhalten werde, auf daß die Unterthanen ein geruhiglich und stilles Leben führen mögen, in aller Gottseligkeit und Ehrbarkeit, Amen. Dies lehret auch das 17. Capitel.

\section[Bestellung, Wahl und Amt der Richter, sonderlich eines]{Bestellung, Wahl und Amt der Richter, sonderlich eines Königs}
\quellseite{304}

\section{Summarium}

Ist eine Ordnung, wie die, so im weltlichen Regiment Richter sind, sich recht halten, und nicht Geschenke nehmen sollen, V. 19., und sonderlich ob dem fest halten sollen, daß keine Abgötterey geduldet werde, V. 21. Was im selbigen Fall mit zweien oder dreien Zeugen kann überwiesen werden, daß man nicht schonen, sondern tödten soll, V. 5.; denn das ist der fürnehmsten Werke eins, daß der Obrigkeit von Gott ist auferlegt, daß sie ob dem Wort Gottes und dem rechten Gottesdienst mit Ernst halten, und allem Aergerniß, so viel möglich ist, mit dem Schwert wehren soll.

\section{Summarium}

Ist ein Befehl, wenn sie zum Könige wählen, und wie sich derselbige halten soll. Da ist sonderlich zu merken, daß er verbeut, der König soll nicht viel Rosse noch Weiber haben, und nicht Gold und Silber sammeln, V. 16.; denn diese drei Stücke sind, damit sich die Obrigkeit gemeiniglich versündiget. Denn viel Rosse, das ist, die große Gewalt, verursachen Frevel und Muthwillen, und machen stolze Herren. Viel Weiber oder Unzucht, machen ruchlose Leute, die beide des Regiments und Gottes nicht achten. Gold und Silber sammeln, geschieht nimmermehr ohne der armen Unterthanen Schaden, so von geizigen Oberherrn übersetzt und ausgesogen werden. Darum verbeut Moses, als ein vernünftiger Regent, solches, und sagt, man solls nicht thun. Das aber soll der König thun, daß er das Wort Gottes stetig lese, und daraus lerne Gott fürchten, V. 18. Welche Obrigkeit solches thut, wird sich für den obgemeldeten dreien Stücken wohlwissen zu hüten.

\section{Votum}

GOtt erhalte unsere liebe Obrigkeit in langem Leben und beständiger Gesundheit,
\quellseite{305}
und helfe, daß sie ferner also regiere, auf daß es Ihm zu Lobe, und den Unterthanen zu Nutz und Wohlfahrt gelange, Amen.

\chapter{Das 18. Kapitel}

\section{Vorrede}

Die Zauberei ist die schöne Tugend, darüber auch der Teufel in der Hölle lachet; und ist wohl zu weisen, mit was Namen dies Laster nur könne genugsam genennet werden, daß ein Mensch, der durch Christum aus des Teufels Reich und Rachen gerissen, auch in der Taufe dem Teufel abgesagt, und Gott dem Herrn zu dienen geschworen hat, nun allererst alle solcher Sünde und Blindes vergessen, sich bloß von Gott abreißen und sich muthwillig und fürsetzlich, dem Teufel zu Dienst, übergeben. Diese überaus greuliche große Sünde läßt uns fliehen, und niemand sage, wie solche Leute den Teufel pflegen fromm zu machen: Eintracht doch alles und eitel gute Worte; denn nimmt ein einfältig, doch deutlich Gleichniß: ein gülden Becher ist ja gut und schön; also ein sammetner Rock auch; aber mit dem Becher Mistpfützen schöpfen, und mit dem sammetnen Rock sich in Koth legen, das ist so viel desto ärgere und größere Schande. Also hier, die Worte sind gut, desto ärger ists, daß sie zur Zauberei, wider Gott, sollen mißbraucht werden. Darauf höret nun an das 18. Kapitel.

\section{Vom Priesterrecht, Abgötterei, und Christo, dem rechten Propheten}

a Hie wird klärlich eine andere Predigt verheißen, denn Moses predigt, welche nicht kann das Gleiche sein, das ganzsam durch Mose gegeben, darum muß es das Evangelium sein. Und dieser Prophet niemand denn Jesu Christus selbst ist, der solche neue Predigt auf Erden gelehrt hat.
\quellseite{306}

b Vermessenheit hie redet Mose von den Propheten, so neue Lehre über die alte und vorige Lehre fürsetzen. Solchen soll man ohne Zeichen nicht glauben. Denn Gott allzeit sein neu Wort mit neuen Zeichen bestätiget. Aber droben Kap. 13, 2. redet er von den Propheten, so wider die alte bestätigte Lehre predigen; diesen soll man nicht glauben, wenn sie gleich Zeichen thun, ut supra.

\section{Summarium}

Hier ist endlich ein Befehl von den Leviten, weil sie des Gottesdienstes sollen warten, daß jedermann sie ihm lassen befohlen sein; darnach warnet er für Abgötterei und sonderlich für Zauberei und Wahrsagern, und sagt klar, wie das vermessene Sünden sind, um welcher willen Land und Leute müssen gestraft und verwüstet werden, V. 10. Am Ende ist eine Verheißung, wie Gott noch einen andern Propheten, denn Mosen, mit einer andern Lehre, denn des Gesetzes, schicken werde, V. 25. Dies ist unser lieber HErr JEsus Christus mit seinem Evangelio, den soll man hören und sein Wort annehmen. Wie es der heilige St. Petrus selbst deutet in den Geschichten der Apostel am 3. Kapitel V. 23.: Wer es nicht thut und andern Lügen-Propheten will gehorchen, deß Seele soll ausgerottet werden.

\section{Votum}

Unser lieber HErr JEsus Christus helfe, daß wir sein Wort hören und auch thun mögen, und regiere uns, daß wir uns für Abgötterei, Zauberei, falscher Lehre und allen andern Sünden allzeit hüten. Amen.

\chapter{Das 19. Kapitel}

\section{Vorrede}

Gott will ernstlich, daß, welcher Menschen-Blut vergießt, daß der wieder von ordentlicher Obrigkeit soll getödtet werden. Daher sagt unser HErr Christus Matth. 26.: Wer das Schwert nimmt, der soll durchs Schwert umkommen. Darum hat Gott auch allezeit mit sonderlichem Ernst über dem fünften Gebot gehalten. Solchen Ernst soll weltliche Obrigkeit wohl bedenken, und derhalben ihrem Amt nach keinen Todtschlag ungerochen noch abbitten lassen. Denn unschuldiges Blut schreiet zu Gott, wie Abels Blut; und wo es an dem Rechtsschuldigen nicht gerochen wird, so muß ein ganz Land solche Sünde tragen, und darum dem Zorn und Strafen Gottes herhalten. Dieß soll auch erinnern das 19. Kapitel.

\section{Ordnung der Freistädte. Falscher Zeugen Strafe}
\quellseite{307}


\section{Summarium}

Moses erholet hie das Gesetz von Freistädten, dahin sie fliehen sollten, so wider ihr Wissen und Willen ohngefähr einen Todtschlag hätten begangen, V. 3.; auf daß sonderlich die weltliche Obrigkeit lerne, wie so viel daran gelegen ist, daß niemand unschuldig umgebracht werde. Und wiederum, daß der Schuldige in seinem Wege soll verschonet werden. Wie der Text hie ernstlich sagt: Deine Augen sollen des Todtschlägers nicht verschonen, sondern du sollst das unschuldige Blut aus Israel thun, daß es dir wohlgehe, V. 13. Wenn nun durch die Obrigkeit Mord und Frevel nicht gestraft werden, daselbst kanns nicht wohl gehen, sondern es müssen sich allerlei gemeine Plagen finden. Denn wie Mose von den falschen Zeugen hernach sagt: Es müssen solche Sünden gestraft werden, auf daß die andern, so es hören, sich fürchten, und nicht mehr solche böse Stücke zu thun fürnehmen.

\section{Votum}

Helfe der treue GOtt, daß wir der Obrigkeit unterthan seyn, und daß die Obrigkeit christlich und wohl regiere, zur Rache über die Uebelthäter, und zu Lobe den Frommen, Amen.

\chapter{Das 20. Kapitel}

\section{Vorrede}

Wie der Land- und Stadtfriede ist ein halber Himmel und ein irdisch Paradeis; also ist der Krieg eine halbe Hölle, oder ja ein Vorbote und Vorläufer der Höllen. Denn wie der Friede eine Gabe Gottes ist, daraus alle andern zeitlichen Güter fließen und quellen: also, wenn Gott mit Krieg straft, so gibt immer ein Unglück dem andern die Hände. Derwegen danken wir Gott heute herzlich, daß wir guten Frieden bisher haben, und laßt uns zuschauen, daß wir also leben und also beten, daß wir den güldenen Frieden auch ferner behalten. Denn die Furcht des HErrn gibt reichen Frieden und Heil, sagt Sirach im 1. Kapitel. Ferner höret davon an das 20. Kapitel.

\section{Kriegsregeln}

a Ein Haus einweihen, heißt auf hebräisch, wenn man drein zieht, drinnen zu wohnen.
\quellseite{308}
Ist's doch Holz auf dem Felde, und nicht Mensch, daß es vor dir ein Bollwerk sein möge. Es ist Holz auf dem Felde, und nicht in der Stadt, es thut dir nichts, und ist dir nütze. Hic sensus congruit allegoria, non esse pugnandum contra eos, qui non sunt contra nos, sed pro nobis. Der Sinn des Gleichnisses ist: man soll nicht gegen die streiten, die nicht gegen uns, sondern für uns sind.

\section{Summarium}

Ist erstlich eine feine Ordnung, wie man dem Kriegsvolk im Streit soll zusprechen, daß sie sich nicht entsetzen sollen für den Feinden, sondern ihre Hoffnung darauf setzen, daß Gott, welcher ihnen aus Egypten und allerlei Noth treulich geholfen habe, bei ihnen sein, für sie mit ihren Feinden streiten und ihnen helfen wolle. Solcher Trost ist gewiß und fehlet nicht, V. 1. Aber es gehöret auch dazu, daß man eine gute Sache habe, denn zum Unrecht will Gott niemand helfen. Welche Obrigkeit aber eine gute Sache hat, die mag Schwert, Büchsen und Harnisch brauchen; aber das Vertrauen soll allein auf Gottes Hülfe stehen, so wird sich der Sieg endlich auch wohl finden.

Darnach sehet Moses etliche andere seine und billige Ordnung, so in Kriegsläuften zu gebrauchen sind, V. 10. Deren zum Theil man sich noch könnte brauchen. Aber es ist leider bei unserm Kriegsvolk seine Zucht noch Gottesfurcht. Moses verbeut, sie sollen keinen fruchtbaren Baum in der Feinde Land beschädigen. Denn, spricht er, es ist ein Holz und nicht ein Mensch, V. 19. So verwüsten unsere Krieger muthwillig auch der Fremden Güter, Weinberge, und verderben, was sie können. Daß man wohl gedenken mag, es sind nicht Menschen, sondern lauter Teufel. Was sollte denn für Glück bei ihnen sein.

\section{Votum}

Der treue, gütige GOtt wende ab die wohlverdienten Strafen, erleuchte die Kriegsleute zu seiner Furcht und Gerechtigkeit, und behüte uns ja für Krieg und Blutvergießen, um JEsu Christi, seines lieben Sohnes und unsers HErrn, willen, Amen.

\chapter{Das 21. Kapitel}

\section{Vorrede}

Es sind nicht allein Mörder die, so mit wehrhafter Faust, oder sonst mit Gewalt und Gift,
\quellseite{309}
einen Menschen todtschlagen und ermorden; sondern die Kinder, welche ihre Eltern übel halten, oder mit ihrem Ungehorsam und Schandleben ängsten, daß sie für Sorge und Grämniß sterben müssen, das sind auch Vater- und Mutter-Mörder. Wehe aber beiden. Denn kein Mörder bleibt ungestraft, und es ist auch nie keinem ungehorsamen Kinde glücklich und wohl gegangen, und wenn sie sich nicht bekehren, so ist ihr letzter Lohn höllisch Feuer und ewige Verdammniß. Denn welche solches thun, werden das Reich Gottes nicht erben, sagt St. Paulus zum Galatern im 5. Kapitel. Dies lehret auch das 21. Kapitel.

\section[Vom unbekannten Todtschlage; gefangenen Weiberpersonen; Recht des Erstgebornen;]{Vom unbekannten Todtschlage; gefangenen Weiberpersonen; Recht des Erstgebornen; ungehorsamen Sohne; Gelübden}
\quellseite{310}

\section{Summarium}

Ist eine Ordnung, wie es erstlich soll gehalten werden, wenn einer im Felde erschlagen wird, da man nicht weiß, wer es gethan hat, daß die Priester und Obrigkeit zusammen kommen, und Gott bitten sollen, daß er solche Sünde nicht am ganzen Volk strafen wolle, weil man den Thäter nicht haben und strafen kann, V. 1. Da sehen wir abermal, wie eine große Sünde es ist, todtschlagen und den Todtschläger nicht strafen. Darnach sind etliche andere Ordnungen, unter welchen sonderlich zu merken ist, wie Gott die ungehorsamen Kinder, so Vater und Mutter übel halten, bestraft, daß man sie tödten soll; auf daß andere sich bessern und lernen fromm und gehorsam sein.

\section{Votum}

Der starke GOtt steure allen blutgierigen Menschen, behüte uns für derselben Zorn und Mord, und regiere die aufwachsende liebe Jugend, daß sie ihren Eltern gehorsam sei und folge, und an Leib und Seele gesegnet werde, Amen.

\chapter{Das 22. Kapitel}

\section{Vorrede}

Unser Seligmacher JESUS Christus sagt Matth. am 5.: Selig sind die Barmherzigen, denn sie werden Barmherzigkeit erlangen. Und: Selig sind die reines Herzens sind, denn sie werden Gott schauen. Daraus folget, daß Gott denen, so unbarmherzig gegen ihrem Nächsten sein, auch keine Barmherzigkeit erzeigen wolle; und die Unzucht lieben und üben, daß sie zeitlich an Leib, Seele, Gut und Ehre sollen gestraft werden, und, da sie in solchen Sünden ohne wahre Buße sterben, daß sie auch in Ewigkeit von Gottes Angesicht, zu ewiger Pein und Marter, verstoßen sein sollen. Dergleichen läßt uns Barmherzigkeit an unserm Nächsten beweisen, und alle Unreinigkeit unsers Leibes fliehen und meiden. Davon höret auch an das 22. Kapitel.

\section{Mancherlei Gesetze verzeichnet}
\quellseite{311}

\section{Summarium}

Ist eine Ordnung, wie man verloren Vieh und anderes bewahren, und dem wieder zustellen soll, der es verloren, V. 1. Sonderlich aber ist hie zu verwundern, daß Gott befiehlt und verheißet: Wer ein Vogelnest ausnimmt, und die Mutter fliegen läßt, daß demselben Menschen wohl gehen und er langes Leben haben soll, V. 6. Denn Gott will, daß wir barmherzig sollen sein, auch gegen das unvernünftige Vieh, und hat ein Ungefallen dran, wo man tyrannisch handelt, wenn es gleich nur um ein armes Vöglein zu thun ist. Wir müssen aber bei diesem Befehl annehmen, so es Gott wohlgefällt, wo man mit den unvernünftigen Thierlein freundlich umgeht, wie viel ein größer Wohlgefallen er an dem haben, und mit reichlichem Segen ers belohnen wolle, wo ein Mensch dem andern freundlich und behülflich ist.

Es ordnet auch Moses hie, wie man Ehebrecher und Jungfrauschänder steinigen und tödten soll, V. 22. Auf daß man lerne, wie große Sünde dies sei, obs wohl weltliche Obrigkeit nicht so heftig strafet, wiewohl vonnöthen wäre. Und ist gar kein Zweifel dran, daß deshalb viel gemeiner Strafen täglich über uns gehen, daß solche Laster so gar keine Strafe haben.

\section{Votum}

Gott breche aller Unbarmherzigen Muth und Sinn, und steure dem Schandgeist und aller Unzucht, und regiere uns und die liebe Jugend, daß wir uns für allerlei Sünde und Schandlastern hüten, und dagegen züchtig und christlich leben, Amen.

\chapter{Das 23. Kapitel}

\section{Vorrede}

Es ist kein größerer Menschenfeind auf Erden nach dem Teufel, denn ein Geizhals und Wucherer; und sollen derwegen fromme Herren solche Wucherer und Geizhälse ansehen als leibhaftige Teufel und sich für ihnen segnen, wo man sie höret oder siehet, und sollen gedenken, daß Türken, Tartarn und andere Heiden eitel Engel sind gegen einem Wucherer. Desgleichen sollte man die Jugend lehren und gewöhnen, daß sie erschrecken und pfui dich sagen für dem Namen Wucherer, als für dem ärgsten Teufel. Weil dem also, so helfe Gott, daß die, welche solcher Sünden sich schuldig wissen, rechte wahre Buße thun, und regiere alle, die an diesem Hauptlaster unschuldig sein, daß sie nimmermehr in diese Sünde, zu ihrem und der Ihren Unheil, gerathen, Amen. Dies erinnert auch das 23. Kapitel.
\quellseite{312}

\section{Wer in die Gemeine des Herrn gehöre, oder nicht}

a Bösen) daß du selbst nicht unsträflich seiest, und also den Sieg zur Strafe verlierst und geschlagen werdest, wie zur Zeit Eli und Saul geschah.

\section{Summarium}

Ist eine Ordnung von mancherlei Sachen, unter welchen sonderlich diese zu merken sind, wie Moses befehlet, wenn man wider die Feinde ziehen soll, daß jedermann sich für dem Bösen hüten soll, v. 9. Denn weil man Gottes Gnade und Hülfe zum Sieg bedarf, soll man sich nicht versündigen und Gott nicht Ursache zur Strafe geben; wie man stehet, wie es mit unserm Kriegsvolk leider gehet, da alle Unzucht und Gotteslästerung
\quellseite{313}
 auf das greulichste geübt wird; darum ist auch so wenig Glück dabei.

Darnach befiehlt er, daß kein Jude am andern soll wuchern, V. 19. Viel mehr sollte solches kein Christ gegen den andern thun. Einmal wir durch Christum zu Brüdern unter einander gemacht und einer dem andern um christlicher Liebe willen, und nicht um Eigennutz willen, gern dienen und helfen sollte, wie denn alle Christen thun. Die es aber nicht thun, sind keine Christen, ob sie gleich den Namen haben; müssen derhalben zum Teufel und in das höllische Feuer.

\section{Votum}

Der starke und eifrige GOtt steure allen, die muthwillig sündigen, und helfe, daß ein jeder gegen seinen Nächsten also handele, daß sich GOtt wieder seiner aus Gnaden erbarme, und ihm helfe an Leib und Seele, bis zum ewigen Leben, Amen.

\chapter{Das 24. Capitel}

\section{Vorrede}

Welcher Mensch sich mit Schaden seines Nächsten nähret, der erarbeitet wahrhaftig nicht, und je glückseliger es ihm am Anfange und eine Zeitlang dazu gehen mag, je mehr Verderben und größerer Schaden endlich folget. Dies lehret mit seiner Vermahnung und Warnung der heilige Apostel Paulus, da er 1 Thess. 4. also sagt: Das ist der Wille Gottes, daß niemand zu weit greife, nicht übervortheile seinen Bruder im Handel; denn der HErr ist ein Rächer über das alles. Derwegen ein jeder seinem Nächsten in seinem Beruf, wo und wie er kann, diene und helfe; er bete auch, und sei gottsfürchtig daneben, so wird Gottes Segen bei seiner Nahrung sein und bleiben. Ferner höret an das 24. Kapitel.

\section{Gesetze, vornehmlich den häuslichen Stand betreffend}
\quellseite{314}

\section{Summarium}

In eine Ordnung von dem Scheidebrief, V. 1., davon Christus klärlich sagt, Moses habe ihn den Juden erlaubet, nicht daß er recht sei, sondern daß sie so ein störrig, unverträglich Volk sind. Darnach ordnet er, wie mans mit Leihen und Pfand nehmen dermaßen halten soll, daß den Armen damit geholfen und sie nicht mehr beschweret werden, V. 10., und sonderlich gebeuget er, man soll die Pfande über Nacht im Hause nicht behalten, der die armen Leutlein ohne ihren Schaden nicht entbehren können, V. 13. Solche und dergleichen andere Lehren, so hie angezeigt sind, sollten wir bedenken und uns darnach halten, so würde der Untreue des Geizes, Schindens und Kratzens weniger werden. Weil es aber nicht geschiehet, sehen wir, daß ein Unglück nach dem andern sich findet.

\section{Votum}

Helffe der treue GOtt, daß sich ein jeder in seinem Beruf treulich und christlich nähre, und seinem Nächsten gern und nutzbarlich diene, damit ihm GOtt wieder diene zu aller Nothdurft, um Christi JEsu willen, Amen.

\chapter{Das 25. Capitel}

\section{Vorrede}

Der heilige Apostel Paulus vergleicht die Geizigen den Abgöttern. Denn gleich wie ein abgöttischer Mensch seinen Schöpfer verläßt, und seine Hoffnung und Zuversicht setzet und gründet auf die Creaturen: also thut auch ein Geizhals, welcher Gold und Silber für seinen Gott hält. Das ist aber sehr schrecklich zu gedenken, will geschweigen zu reden, daß ein Mensch eine stumme tote Kreatur seinem Schöpfer soll vorziehen. Zu dem, so ist das Gottes Ordnung für und für, wo er das Leben gibt, daß er auch schafft, daß man es erhalten könne. Dieweil er denn solches dem Viehe thut, wie viel mehr will er solches den Menschen und sonderlich seinen Christen thun, welchen er nicht allein Leib und Leben, sondern auch seinen eingebornen Sohn gegeben hat, daß sie nicht allein zeitlich, sondern in Ewigkeit leben sollen. Dergleichen für dem Geiz sich jedermann hüten soll. Darauf höret auch davon an das 25. Kapitel.

\section{Etliche Stücke jüdischer Polizeiverordnung}
\quellseite{315}

\section{Summarium}

Ist erstlich eine Ordnung, wie man die mit Schlägen strafen soll, so bußfällig und dennoch den Tod nicht verdienet hätten. Denn Sünde soll nicht ohne Strafe hingehen. Darnach ordnet er anders mehr, und ist sonderlich zu merken, wie er bräuet, wer unrecht Gewicht, Maß, Ellen brauchet, daß ihn Gott strafen werde. Wer aber recht Gewicht und gute völlige Maße gebe, den werde Gott segnen, V. 15. Solche Sprüche sollten die Kaufleute in die Kräme an Waage und Gewicht schreiben und malen lassen, und sich darnach halten. Aber der leidige Geiz stopfet ihnen die Ohren, und hält ihnen die Augen zu, daß sie es verachten. Darum verursachen sie Gott, daß es unter ihnen wüste umgehet, jetzt da einen, jetzt dort einen am Leib, Gut und andern solcher Sünden halben strafen muß. Es wird aber endlich, wo Besserung nicht folget, auch das höllische Feuer über sie kommen, wie Paulus sagt: Die Geizigen werden das Reich Gottes nicht erben.

\section{Votum}

Verleihe der treue GOtt, daß die Obrigkeit die Sünden strafe, und daß für Betrug, Geiz, und anderen Sünden wir uns hüten, und daß wir dem Nächsten brüderlich rathen und dienen, um unsers lieben Bruders Christi Jesu willen, Amen.

\chapter{Das 26. Capitel}

\section{Vorrede}

Die heilige Schrift lehret, daß die Kinder Gottes sein Erstlinge seiner Kreaturen, Jacobi 1, 18. (Gleich wie man aber im alten Testament die Erstlinge allerlei Früchte dem HErrn geheiliget und ergeben, oder Gott zur Dankbarkeit geopfert hat: also sollen wir uns auch im neuen Testamente Gott heiligen und opfern als Erstlinge seiner Kreaturen und seine neugborne Kinder, das ist, wir sollen unsern Leib und alle innerliche und äußerliche Gliedmaßen also gebrauchen, daß Gott täglich von uns geehret und gehorsamlich dienet werde, und daß wir nach seinem Willen leben, ein jeder nach seinem Beruf, Amt und Stande. Davon lehret St. Paulus Röm. 12, 1. also: Euer Leiber gebet zum Opfer, das da lebendig, heilig und Gott wohlgefällig sei, welches sei euer vernünftiger Gottesdienst. Dies erinnert auch das 26. Kapitel.

\section{Von den ersten Früchten und Zehnten}
\quellseite{316}

a Leide) Gottes Opfer soll fröhlich, heilig und rein sein; darum soll nichts in Traurigkeit davon gegessen, nichts in Unreinigkeit davon genommen, nichts den Götzen oder Todten davon gegeben sein.

\section{Summarium}

Dieß ist die Ordnung von den Erstlingen der Früchte, von Zehnten, V. 2., in welcher wir sehen, wie Gott immerdar befiehlt, daß man seinen Segen mit Dank annehmen, und denselbigen dazu brauchen soll, daß zum ersten die Leviten, das ist die, so die Kirchendämter verwalten, versorget, und darnach Wittwen, Waisen und Fremdlingen damit geholfen werde. Wer nun sein Gut also brauchet, der mag sich trösten, wie die Juden hie bitten, daß GOtt solches sehen, und sie und das Land mildiglich segnen werde, V. 12., da sonst bei den Geizigen und die sich der Dürftigen nicht annehmen, weder Glück noch Heil ist.

\section{Votum}

Der treue GOtt steure allem Geiz und Unbarmherzigkeit; und die des Armuths und seiner Kirchen sich treulich annehmen, die segne GOtt an Leib und Seele, zeitlich und ewig, um JEsu Christi seines lieben Sohns willen, Amen.

\chapter{Das 27. Capitel}

\section{Vorrede}

Was Gott den Unbußfertigen dräuet in seinem Wort, das läßt er wahrhaftig ergehen und geschehen, wie lange auch die Strafe außen bleibt. Derwegen, ach, laßt uns des heiligen Wortes Gottes Flüche und Dräuungen nicht mit tauben Ohren und steinernen Herzen, sondern mit Andacht fleißig anhören; und erschrecket alle für
\quellseite{317}
Gottes Zorn, hütet euch für allerlei Sünden, und lebet, wie es euch Gott ernstlich in seinem Gesetz und Wort befohlen hat. So wird nicht Gottes Fluch, sondern sein Segen über euch kommen, und also wird Gott euer lieber Vater, ihr seine Kinder, der HErr Christus euer treuer Bruder, ihr auch seine Brüder und Schwestern, der heilige Geist euer beständiger Lehrer und Tröster, ihr seine Tempel und Wohnungen sein und bleiben. Denn also sagt der HErr Christus auch selber Joh. 14, 23.: Wer mich liebet, der wird mein Wort halten, und mein Vater wird ihn lieben, und wir werden zu ihm kommen und Wohnung bei ihm machen. Darauf höret nun das 27. Kapitel.

\section{Denksteine der Kinder Israel. Segen und Fluch}

\section{Summarium}

Ist eine Ordnung, wie die Leviten öffentlich verfluchen sollen, alle, die wider das Gesetzes
\quellseite{318}


HErrn thun, und die ganze Gemeine soll dabei stehen und Amen dazu sagen, V. 15. In solchem Fluch werden unter andern sonderlich genennet, ungehorsame Kinder, so ihren Eltern fluchen, V. 16.; auch weltliche Obrigkeit, die den Fremdlingen, Wittwen und Waisen nicht zum Rechten verhelfen, sondern Unrecht thun, V. 19. Und ist hie sonderlich zu merken, daß solcher Fluch allein auf die Sünden gehe, so ohne Mittel wider Gott und das Gewissen, oder wie mans sonst nennet, wider das natürliche Recht sind.

\section{Votum}

Der Vater aller Barmherzigkeit, GOtt, unser lieber HErr, sei uns allen gnädig, helfe, daß wir in seinen Geboten wandeln, und erhalte uns in seinem Segen, an Leib und Seele, bis zum ewigen freudenreichen Leben, das uns erworben hat unser HErr JEsus Christus, Amen.

\chapter{Das 28. Capitel}

\section{Vorrede}

Das Wort Gottes ist der unvergängliche Same, wie St. Petrus spricht, dadurch wir neu geboren werden. Sintemal der Glaube, der uns zu Kindern Gottes macht, kömmt aus der Predigt des Evangelii, dadurch wir in unsern Herzen, Verstand und Willen erneuert werden, daß, wie wir vorhin unserer ersten Geburt nach, ohne alle Erkenntniß und ohne alle Furcht Gottes lebten und unser Herz voll böser Lust und Begierde waren: also gewinnen wir, der andern Geburt nach, ein herzliches Mißfallen über die Sünde, und ein Gefallen, daß wir Gott fürchten und lieben, ihm vertrauen, daß wir uns befleißen eines gottseligen Lebens und Wandels, und daß wir Lust und Freude am Gesetz des HErrn haben. So helfe Gott, daß zu unserer christlichen Besserung sonderlich viel diene das 28. Kapitel des fünften Buches Mos.

\section{Verheißener Segen, gedrohter Fluch}

a Korb) das ist alles, was du gegenwärtiglich brauchest, und was du beiseit legest zu behalten.

\section{Summarium}

Es ist die Beschluß-Predigt. Denn nachdem Moses wiederholet hat die fürnehmsten Stücke, was sie thun und lassen sollen, wenn sie Gott gefallen und recht thun wollen, beschleußt er

18*
\quellseite{319}
nun und sagt erstlich, wie Gott sie segnen und solchen Gehorsam gegen seinem Wort reichlich bezahlen wolle, mit allerlei Gaben und Gnaden, V. 1. Auf daß wir auch lernen, wer da begehrt, daß er Glück und Heil haben möge, daß er sich nach Gottes Wort und Willen halte. Wie Christus sagt Matth. 6, 33.: Trachtet am ersten nach dem Reiche Gottes und seiner Gerechtigkeit, so wird euch das andere alles zufallen.
\quellseite{320}

b Alii infantes recens natos. D. i. Andere: neugeborne Kinder.

\section{Summarium}

Ist das andere Stück der Bußpredigt Mose, in welcher er warnet, daß man sich vor Sünden hüten wolle, sintemal Gott dieselbigen mit allerlei Plagen strafen will. Da sollt du lernen, was die rechte Ursache sei, daß Krieg, Pestilenz und andere Krankheit, Hunger und Kummer in der Welt ist, nämlich, daß solches unsere eigene Sünden verursachen. Darum sollten wir von Sünden ablassen, so würden solche Strafen auch ausbleiben.
\quellseite{321}

\section{Summarium}

Dies ist eine helle Anzeigung, warum solche grausame Plagen kommen sollen, nämlich, daß man der Stimme des HErrn nicht gehorchet, und nicht thut, was Gott in seinem Gesetz befohlen hat. Und ist sonderlich hie zu merken, daß Gott an solchen leiblichen Strafen nicht gesättiget ist, sondern auch die höchsten und geistlichen Strafen dräuet, daß man andern Göttern, Holz und Stein, mit unruhigem, erschrockenen Gewissen dienen soll, wo man dem rechten Gott nicht dienen und seinem Wort nicht hat mit gutem Gewissen folgen wollen, v. 64.

\section{Votum}

Verleihe der treue, barmherzige GOtt, daß wir dieser Lehren nimmer vergessen, sondern daß wir nach diesen allezeit leben, in Gerechtigkeit und Heiligkeit, und also thun was recht ist, daß uns GOtt Gutes thue, wie er uns verheißen hat, Amen.

\chapter{Das 29. Capitel}

\section{Vorrede}

Wer Gottes gutes heiliges Wort höret, und doch böse und sündliche Werke thut, den bedrohet Gott mit Bösem. Denn wie sich ein Mensch gegen Gottes Wort hält, also hält sich Gott wiederum gegen denselbigen Menschen. Derwegen wer nicht will, daß ihn Gott verachten soll, der verachte nicht Gottes Wort. Und wer da will, daß ihm Glück und Segen folgen soll, der folge auch Gottes Wort, daß er nach demselbigen gehorsamlich lebe. Denn die Gottes Wort hören und thun, oder wer den Willen thut Gottes, des Vaters im Himmel, derselbige ist des HErrn Christi, Matth. 12, 50., das ist, er will ihnen Liebe und Treue erzeigen zu aller Wohlfahrt, mehr denn das Geschwister einander, ja mehr, als die Eltern ihren Kindern Gunst und Dienst erzeigen können. Darauf höret ferner an das 29. Kapitel.

\section{Erneuerung des Bundes}
\quellseite{322}

a Eurem Thun) ohne Gottes Wort ist all unser Thun Narrheit.

b Es gehet mir wohl) das ist der ruchlosen Leute Wort und Gedanken: ei die Hölle ist nicht so heiß; es hat nicht Noth, der Teufel ist nicht so greulich, als man ihn malet. Welches alle Werckheiligen frech und durstiglich thun, ja noch Lohn im Himmel gewarten.

c Weil ich wandele) das ist, weil ich so thue, und meine Abgötterei treibe, so ist eitel Glück da; wer Gott dienet, dem gehets nimmermehr wohl, Apost. 8, 23.
\quellseite{323}

d Das Geheimniß will so sagen: Uns Juden hat Gott für allen Völkern auf Erden seinen Willen geoffenbaret, und was er im Sinn hat. Darum sollen wir auch desto fleißiger sein.

\section{Summarium}

Das ist eine Verneuerung des Bundes, weil sie alle Wunderwerke in Egypten gesehen, V. 2., und hernach erfahren haben, wie sie Gott in der Wüste erhalten, V. 5., ihnen wider die Heiden Sieg gegeben, und sie in das verheißene Land eingesetzt hat, daß sie sich Gottes Willens und seines Wortes halten sollen, und sonst keinem andern Gott dienen, sonst soll alles Unglück über sie kommen, V. 18. Solches soll uns eine Verwarnung sein, daß wir uns für Sünden hüten und nach Gottes Wort thun sollen, sonst wird dergleichen Fluch über uns auch kommen.

\section{Votum}

GOtt, unser treuer, lieber Vater, lasse seinen Segen über uns und all den Unsern bleiben, und helfe, daß wir nach seinem Wort thun und leben, wie er uns befohlen hat, und wir ihm solches schuldig sind als seine Kinder, Amen.

\chapter{Das 30. Kapitel}

\section{Vorrede}

Gottes Barmherzigkeit hat noch kein Ende, sondern sie ist alle Morgen neu, sagt der heilige Prophet Jeremias in Klagliedern im 3. Kap. V. 22. Und die heilige Jungfrau Maria singet schön und tröstlich in ihrem Magnificat, Luc. 1, 50.: Gottes Barmherzigkeit währet immer für und für, bei denen, die ihn fürchten. Diese Barmherzigkeit Gottes sollen wir also gebrauchen, daß ein jedes seine Sünde erkennen, Gott dieselbigen alle herzlich abbitten, und sich trösten soll, obschon seiner Sünden viel und die größten sind, daß Gottes Barmherzigkeit tausendmal mehr und größer ist, der ihm alle seine Sünde vergeben will; doch, daß niemand auf Gottes Barmherzigkeit muthwillig sündige. Denn wer auf Barmherzigkeit sündiget, dem wird mit Unbarmherzigkeit endlich abgelohnt. Darauf höret an das 30. Kapitel.

\section[Die Bußfertigen erlangen Gnade, die Widerspenstigen Rache und]{Die Bußfertigen erlangen Gnade, die Widerspenstigen Rache und Strafe}

a Zu gut denn die Gottlosen haben auch wohl Ehre und Gut, oft mehr, denn die Heiligen, aber zu ihrem und Andern Verderben.
\quellseite{324}

\section{Summarium}

Ist eine tröstliche Verheißung, wie Gott wiederum zu Gnaden sein Volk wolle annehmen, wenn sie nur in sich selbst schlagen und lassen ihnen ihre Missethat leid sein, und begehren sich zu bessern, V. 2. Solche Verheißung gehet fürnehmlich auf das Neue Testament. Obwohl das Volk ihrer Sünden halben mancherlei soll gestraft werden, daß dennoch Gott Christum senden, und durch den heil. Geist ihre Herzen beschneiden und sie recht fromm machen werde. Darum vermahnet er am Ende, weil er beides ihnen habe fürgelegt, das Gute, so denen nachfolgen soll, die Gott lieben, und sein Wort halten, und den Fluch, darauf die Sünden gehet, daß sie sich zum Guten halten und fromm sein sollen.

\section{Votum}

Der treue, gütige GOtt sei uns allen gnädig, und vergebe uns alle unsere Sünde, und strafe uns ja nicht nach Verdienst, sondern segne und behüte uns nach seiner großen Barmherzigkeit, um Jesu Christi, unsers Heilandes und Mittlers, willen, Amen.

\chapter{Das 31. Capitel}

\section{Vorrede}

Gott schicket uns Kreuz und Unglück zu, daß er den zukünftigen Sünden wehre und uns des Fleisches Kitzel vertreibe. Denn wenn wir ohne Kreuz und Widerwärtigkeit sein, und es uns nach unserm Wunsch und Willen gehet, so wird unser Fleisch fürwitzig, muthwillig und sicher, und gewinnet Lust und Liebe zu sündhaften Dingen; aber wenn wir unter dem Kreuze liegen, so vergehet uns die Geilheit. Daß also gewiß ist, wir sind frömmer, wenn es uns übel gehet. Und dies beweiset das Exempel des Königs Davids. Wie heilig war doch dieser Herr, da ihn Saul und darnach Absalon verfolgte; aber da er Ruhe und Friede hatte, fiel er und sündigte schändlich. Derwegen wir die Vater-Ruthe unsers lieben Gottes küssen und in Geduld das Kreuz tragen sollen. Davon höret auch an das 31. Capitel.

\section{Mose sagt sein Amt auf; ordnet Josua an seine Stelle}
\quellseite{325}

\section{Summarium}

Sie ist ein fein Exempel, wie Moses erstlich das Volk, darnach seinen Diener Josua, der fortan ihr Hauptmann sein sollte, tröstet, daß sie beide auf GOtt ihr Vertrauen setzen, und sich für den Feinden nicht entsetzen sollen, V. 3.

Denn wer GOtt auf seiner Seite hat, der wird auch wohl mitten im Tode sicher sein; dagegen wo GOtt nicht dabei ist, und die Hand abzeucht, da kann weder Geld noch Gut, weder Gewalt noch großes Vermögen helfen oder retten. Solchen Trost soll alle Obrigkeit fürnehmlich haben und durch fleißiges Gebet täglich suchen.
\quellseite{326}

\section{Summarium}

Ist ein fein Exempel, wie Gott gern aller Ursache vorkommen wollte, daß sein Volk nicht sündigte. Aber Art läßt von Art nicht, sondern es gehet, wie der HErr hie sagt: Wenn wir wohl stecken und genug haben, so vergessen wir unsers HErrn Gottes und seines Worts. Darum muß er mit allerlei Unglück hinter uns her sein, auf daß wir nicht sicher werden und in Sünde fallen.

Eben um dieser Ursache willen befiehlt er droben, daß man dies Buch für dem ganzen Volk lese, und das Lied, so hernach folget, singen soll, V. 24. Denn wo uns Gottes Wort nicht erhält, ists nicht möglich, daß wir anders können, denn sicher und gottlos sein, wenn es uns wohl gehet und keinen Anstoß haben.

\section{Votum}

Der treue Gott helfe, daß wir nach seinem Wort in Gehorsam leben, ihm vertrauen, und unter allem Kreuz von ihm Trost, Geduld und Hülfe haben und behalten, bis zur fröhlichen Erlösung, Amen.

\chapter{Das 32. Kapitel}

\section{Vorrede}

Weil die Kinder sind die edlen Pflänzlein, welche Gott in seinem Lustgarten des Ehestandes sollen erzogen werden; so sollen die Eltern nicht gedenken, es sei genug und alles gethan, wenn die Kinder getauft sind. Es ist wahr, sie sind als eine wilde Art abgebrochen und durch die heilige Taufe eingepfropfet und einverleibet der Unschuld Jesu Christi. Aber da sollen die Eltern thun, als fleißige Gärtner und Haushalter des HErrn. Wie ein milder Stamm noch immer ausschlägt in Wasserreiser, und der Gärtner dem Bösen kaum wehren, die guten Reiser aber schwerlich fortbringen kann: also gehet es auch mit den Kindern. Die alte Adams-Art will immer herfür, wie die Schobreiser. Soll nun das Gute, das in sie gepflanzet, nicht zum Abnehmen kommen, so müssen die Eltern ihnen wehren, und mit christlicher Zucht die Uebermaße abschneiden, und solches, weil der Baum noch zart und jung ist. Ach dieses helfe allen lieben Eltern der treue Gott, Amen. Darauf höret auch an das 32. Kapitel.

\section{Moses Lied und Lobgesang}

a Allein) laßt unserm Gott allein Gott sein, und alle Ehre haben und keinem andern.

b Fels) die hebräische Sprache heißt Gott einen Fels, das ist, ein Trotz, Trost, Hort und Sicherung, allein, die sich auf ihn verlassen und ihm vertrauen.

c Gerecht) bei den Gottlosen muß Gott immer unrecht haben und sich meistern lassen, Matth. 11, 19.: Die Weisheit muß sich lassen rechtfertigen von ihren Kindern, die wissen alles besser; wie es Gott macht, so taugts nicht.
\quellseite{327}

d Nicht kannten) hatten sie keinen Befehl noch Gottes Wort gehabt, sondern solchen neuen Gottesdienst aus eigener Andacht erwählet haben.

\section{Summarium}

Sie ist ein Lied, in welchem Moses das Volk lehret, wie Gott ihnen so gnädig gewesen und so große Hülfe und Wohlthat ihnen bewiesen hat, daß ers zum eigenen Volk angenommen, aus Egypten geführet, ihnen das Gesetz gegeben, und endlich ins gelobte Land gebracht hat, V. 7. Darum vermahnet er nach dem ersten Gebot, daß sie solchen Gott ehren, auf seine Hülfe allein sich verlassen und sonst keinem fremden Gott dienen sollen. Aber er weissaget, daß sie bei solcher Lehre und Pflicht nicht bleiben werden; sondern wenn sie satt, dick und fett sein werden, werden sie solchen Gott fahren lassen und Abgötterei treiben, V. 15. Und zum Exempel, daß wir nimmer fleißiger beten und weniger sicher sein sollen, denn wenn wir Friede haben und es uns wohl geht.

e Beräuben, schrecken) das ist: außen sollen sie Wittwen und Waisen werden durch Schwerdt und Gefängniß der Kinder, Männer und Weiber beraubt. Was aber innen bleibt, soll durchs Schrecken, das ist, durch Hunger, Pestilenz, Aufruhr jämmerlich umkommen.

f Rath) sie achten Gottes Wort nicht, wissen alles besser.

g Versiegelt) sie glauben nicht, bis sie es erfahren. Denn es ist für ihren Augen verborgen.
\quellseite{328}

b Custoditum et neglectum, id est, servatam rem et contemptam, quia tempore pacis contemptum est, quod tempore male libenter colligeres, si adesset pretiosum et vile. Das ist, das Verwahrte und Verwahrlosete, d. i. das Beobachtete und Verachtete; weil zur Zeit des Friedens und Wohlergehens verachtet ist, was man zur Zeit des Unglücks und der Noth gern zusammenhielte, wenn es da wäre, Kostbares und Geringes.

i Ueber dem Blut) das sind drei Strafen des Schwerts. Die erste, daß ihrer viel erschlagen werden. Die andere, daß sie gefangen geführet werden. Die dritte, daß ihr Haupt soll bloß werden, das ist, ihr Königreich und Priesterthum soll von ihnen genommen werden, welches durchs Haar auf dem Haupt bedeutet ist.

\section{Summarium}

Ist der andere Theil dieses Liedes, in welchem Moses weissaget, wie sie um Abgötterei und anderer Sünde willen sollten mit allerlei Jammer und Unglück geplaget und endlich ganz und gar verstöret und vertilget werden, und Gott ihm ein ander Volk, die Heiden, werde auserwählen, wie denn geschehen ist. Alles darum, daß sie Gott und seines Worts nicht achteten. Uns zum Exempel. Denn so Gott seines eigenen Volks nicht hat verschonet, dürfen wir nicht denken, wenn wir in gleiche Sünde fallen, daß wir nicht auch gleichen Zorn und Strafe tragen sollen.

\section{Summarium}

Es ist ein fein Exempel, wie Mose sonderlich den Eltern befiehlet, daß sie ihre Kinder zu Gottes Wort halten, ihnen dasselbe für sagen und mit Ernst zu halten befehlen sollen. Denn es sei ein Wort des Lebens, sintemal die es verachten und nicht mit Fleiß halten, darum sterben und verderben müssen. Solche Lehre ist trefflich hoch vonnöthen. Denn die Furcht Gottes, so allein aus dem Wort Gottes kömmt, ist der beste Zuchtmeister; wo dieselbe nicht ist, werden die Kinder ungerathen, und leben ihren Eltern zu Schanden und Schaden. Solchem Jammer ist nicht besser zu wehren, denn man hebe in der Zeit an, die Kinder zur Kirche und Gottes Wort zu gewöhnen, wie Moses hie heißt. Darum ist ein Exempel, wie Moses seiner Sünden halber gestraft wird, daß er nicht in das gelobte
\quellseite{329}
Land kommen soll, V. 48. Wie der Psalm solches auch anzeigt, Gott habe die Sünde seiner Heiligen gestraft; auf daß wir Gott und seinen Zorn lernen fürchten, und uns mit Ernst für Sünden hüten.

\section{Votum}

Regiere uns alle der barmherzige GOtt, daß wir für Sünden uns hüten, und in wahrer Buße leben und sterben. Derselbe treue GOtt helfe auch, daß alle Eltern ihre Kinder wohl ziehen, ihm zu Lob, ihnen selbst zu Ehren, und zur Wohlfahrt den Kindern, Amen.

\chapter{Das 33. Kapitel}

\section{Vorrede}

Ein christlicher Wunsch ist ein christliches und kräftiges Gebet; denn solche herzliche Wünsche sind gegründet auf die Zusagung Christi: Wahrlich, wahrlich, ich sage euch, was ihr den Vater bitten werdet in meinem Namen, das wird er euch geben, Joh. 16, 23. Und wenn man solchen Wunsch und Gruß im Glauben zu allem Dank, mit Freuden aufnimmt, so können sie uns nicht fehlen; denn eines Christen Mund redet nicht vergebene tote und unkräftige Worte, sondern Gott wirket dadurch. Aber wie ihrer viel jetzt einander grüßen, und was sie einander wünschen, das erfähret man oft in betrübten Exempeln. Derwegen wie wir uns selbst alles Gute gönnen, und von Gott bitten, also sollen wir unsern Nächsten auch lieben, und sollen ihm alle Wohlfahrt gönnen, und daß sie ihm widerfahre, wünschen und Gott darum bitten. Dies soll auch erinnern das 33. Kapitel.

\section{Moses Baltesegen und Weissagung}

a Amt) er war nicht König, hatte auch nichts davon, und hielt doch das Volk zusammen, also, daß es ein Haupt hatte, die einen König, und nicht zerstreuet in der Irre ginge.

b Recht) das ist, wie 2 Mos. 28, 30. stehet, daß Heiligthum auf dem Brustlappen. Will also sagen: dein priesterlich Amt sei glücklich für Gott und den Menschen, mit Beten und Lehren.

c Liebliche) das ist, der Tempel und Jerusalem und Königreich war in Benjamin.

d Edle Früchte) das ist vom Königreich Israel gesagt, welches hoch gesegnet ward, mit allem, das Himmel, Sonne, Mond, Erde, Berg, Thal, Wasser und alles zeitliche Gut trug und gab, dazu auch Propheten und heilige Agenten hatte.
\quellseite{330}
• Auszugs) dies hat Debora und Barak gethan, Richt. 5, 12.

* Den Segen Gad hat der König Jehu aufgerichtet, 2 Kön. 10, 27., da er Baal vertilget und das Volk weiter zurechtbrachte, und schlug zween Könige todt, dazu auch Jesebel. Ihr Lehrer ist Elia, der Prophet, der in Himmel genommen und verborgen; denn er war ein Bürger aus Gilead, im Stamme Gad.

§ Gottes Wohnung ist im Volk Israel von Anfang, unter daß sie bleiben ewiglich; als die Hüttlein unter den Armen oder Flügeln der Cherubim.

\section{Summarium}

Ist eine letzte Predigt, in welcher Mosis sein Volk, die zwölf Stämme einen nach dem andern, für seinem Ende segnet, Glück und Heil wünscht, wie im 1 Buch Mos. 49. der Patriarch Jakob auch gethan hat. Aber der Segen hie gehet fürnehmlich auf das zeitliche Priestertum und Königreich, und den zeitlichen Segen, so bei den Juden bis auf Christum, nach Gelegenheit der Zeit, bleiben sollte; wie solches die Historia nach einander fein anzeiget.

\section{Votum}

Der Sohn GOttes, unser Hoherpriester, wolle uns auch segnen, geistlich und leiblich, und geben, was wir sammt den Unsern zeitlich und ewig bedürfen, und wolle solches thun um der Ehre seines heiligen Namens willen, Amen.

\chapter{Das 34. Kapitel}

\section{Vorrede}

Wie andere vor uns gelebt haben und gestorben sind, also die wir jetzund leben, werden auch sterben müssen. Wir leben aber und wissen nicht wie lange, also sterben wir, und wissen nicht wann. Derwegen so sollen wir uns auf ein selig Todtsündlein bereiten, und alle Tage in christlicher Buße leben, auch GOtt bitten, daß er in unserer letzten Noth uns auch beistehe, tröste und stärke wider alle Anfechtungen, und daß er uns erhalten wolle in wahrem Glauben, mit Leib und Seele, zum ewigen Leben; daß uns also der zeitliche Tod eine Thür und Pforte sei, dadurch wir mögen eingehen, aus des HErrn Jesu Christi Gnadenreich dieser Welt, in sein Reich der Glorie und Herrlichkeit, zu ewiger Freude und Seligkeit, Amen. Darauf höret ferner an das 34. und letzte Kapitel des fünften Buchs Mose.

\section{Moses Tod, und was darauf erfolget}
